\documentclass[11pt]{article}
\usepackage[a4paper,margin=22mm]{geometry}
\usepackage{fontspec}
\setmainfont{lmroman10-regular.otf}[BoldFont=lmroman10-bold.otf,ItalicFont=lmroman10-italic.otf,BoldItalicFont=lmroman10-bolditalic.otf]
\usepackage{amsmath,amssymb,amsthm,mathtools,graphicx,xcolor,enumitem,hyperref,xurl}
\setlength{\emergencystretch}{4em}
\newtheorem{theo}{Theorem}[section]
\newtheorem{lemm}[theo]{Lemma}
\newtheorem{prop}[theo]{Proposition}
\newtheorem{coro}[theo]{Corollary}
\newtheorem{corol}[theo]{Corollary}
\theoremstyle{definition}
\newtheorem{defi}[theo]{Definition}
\newtheorem{exam}[theo]{Example}
\newtheorem{exem}[theo]{Example}
\newtheorem{rema}[theo]{Remark}
\newtheorem*{rema*}{Remark}
\newtheorem*{defi*}{Definition}
\newtheorem*{theo*}{Theorem}
\newtheorem*{lemm*}{Lemma}
\newtheorem*{prop*}{Proposition}
\newenvironment{enonce*}[1]{\par\medskip\noindent\textbf{#1.}\itshape}{\par\medskip}
\newenvironment{proofc}{\begin{proof}}{\end{proof}}
\newcommand{\Mc}{\mathcal{M}}

\numberwithin{equation}{section}
%~Mouliné par MaN_auto v.0.25.0 2021-11-08 10:18:54











\renewcommand{\a}{\alpha}
\renewcommand{\b}{\beta}
\newcommand{\g}{\gamma}
\newcommand{\G}{\Gamma}
\renewcommand{\d}{\delta}
\newcommand{\D}{\Delta}
\renewcommand{\i}{\iota}
\renewcommand{\l}{\lambda}
\newcommand{\m}{\mu}
\renewcommand{\r}{\rho}
\newcommand{\s}{\sigma}
\renewcommand{\SS}{\Sigma}
\renewcommand{\t}{\tau}
\newcommand{\e}{\varepsilon}
\renewcommand{\epsilon}{\varepsilon}
\newcommand{\f}{\varphi}
\newcommand{\F}{\Phi}
\newcommand{\p}{\psi}




\newcommand{\Cc}{\mathcal{C}}
\renewcommand{\Mc}{\mathcal{M}}
\newcommand{\Ec}{\mathcal{E}}
\newcommand{\Hc}{\mathcal{H}}
\newcommand{\Tc}{\mathcal{T}}
\newcommand{\Gc}{\mathcal{G}}
\newcommand{\Pc}{\mathcal{P}}
\newcommand*\Uc{\mathcal{U}}
\newcommand*\Vc{\mathcal{V}}
\newcommand{\Wc}{\mathcal{W}}


\newcommand{\R}{\mathbb{R}}
\newcommand{\T}{\mathbb{T}}
\renewcommand{\H}{\mathbb{H}}


\newcommand*\id{\mathrm{id}}
\newcommand*\inj{\mathrm{inj}}
\newcommand*\Aut{\mathrm{Aut}}
\newcommand*\Mod{\mathrm{Mod}}
\newcommand*\Homeo{\mathrm{Homeo}}
\newcommand*\tr{\mathrm{tr}}
\newcommand*\reg{\mathrm{reg}}
\newcommand*\Klein{\mathrm{Klein}}
\newcommand*\geo{\mathrm{geo}}
\newcommand*\Hess{\mathrm{Hess}}
\newcommand*\Sub{\mathrm{Sub}}
\newcommand*\rot{\mathrm{rot}}
\newcommand*\inv{\mathrm{inv}}
\newcommand*\ex{\mathrm{ex}}
\newcommand*\Met{\mathrm{Met}} 

\DeclareMathOperator{\Diff}{Diff}
\DeclareMathOperator{\Rel}{Re}

\newcommand*\wt{\widetilde}
\newcommand*\wbar{\overline}
\newcommand*\wh{\widehat}

\renewcommand*\p{\partial}










%%%%%%%%%%%%%%%%%%%%%%%%%%%%%%%%%%%%%%%%%%%%%%%%%%%
\graphicspath{{./figures/}}

\newcommand*{\mk}{\mkern -1mu}
\newcommand*{\Mk}{\mkern -2mu}
\newcommand*{\mK}{\mkern 1mu}
\newcommand*{\MK}{\mkern 2mu}

%\hypersetup{urlcolor=purple, linkcolor=blue, citecolor=red}

\newcommand*{\romanenumi}{\renewcommand*{\theenumi}{\roman{enumi}}}
\newcommand*{\Romanenumi}{\renewcommand*{\theenumi}{\Roman{enumi}}}
\newcommand*{\alphenumi}{\renewcommand*{\theenumi}{\alph{enumi}}}
\newcommand*{\Alphenumi}{\renewcommand*{\theenumi}{\Alph{enumi}}}
\let\oldtilde\tilde
\renewcommand*{\tilde}[1]{\mathchoice{\widetilde{#1}}{\widetilde{#1}}{\oldtilde{#1}}{\oldtilde{#1}}}
\let\oldhat\hat
\renewcommand*{\hat}[1]{\mathchoice{\widehat{#1}}{\widehat{#1}}{\oldhat{#1}}{\oldhat{#1}}}
%%%%%%%%%%%%%%%%%%%%%%%%%%%%%%%%%%%%%%%%%%%%%%%%

\begin{document}
\section*{2553: Unique least-energy hyperbolic metrics for connected finite weighted graphs surjecting onto the fundamental group of a closed surface of genus greater than one}
\noindent\textbf{Author:} Toru Kajigaya; Ryokichi Tanaka\par\noindent\textbf{Source:} Uniformizing surfaces via discrete harmonic maps. Annales Henri Lebesgue 4 (2021), publisher version, DOI 10.5802/ahl.116. \url{https://ahl.centre-mersenne.org/item/AHL_2021__4__1767_0/}
\par\noindent\textbf{License:} CC BY 4.0, \url{https://creativecommons.org/licenses/by/4.0/}. The original publisher PDF has the explicit license notice. See the saved source/license records.
\par\noindent\textbf{Editorial material note (not author text):} The selected problem is Theorem 2.6. Retained author text comprises the introductory setup, Sections 2--4 and Appendices A--B, including complete local Hessian, convexity, properness, variation and smooth-dependence proofs. The unrelated finite-symmetry section and illustrated examples are omitted; no necessary core figure occurs in the retained proof. All proof text and formulas below are editable author TeX. Mathematical wording, language and proof order are retained. Only the article class, title matter and bibliography presentation are adapted for independent compilation; no proof has been generated or repaired.
\par\noindent\textbf{Editorial difficulty:} H2: The authored argument combines a discrete harmonic-map Hessian along Weil–Petersson geodesics, graph-vertex boundary cancellation and a holomorphic quadratic-differential positivity argument with a collar/Mumford compactness construction and stabilization of markings through surjective fundamental groups. The local new geometric proof is substantially beyond ordinary qualification material.
\medskip\hrule\medskip
\noindent\textbf{Author statement, complete proof and local supporting text follow in source order.}
\section{Introduction}


The Koebe--Andreev--Thurston theorem states that any skeleton of triangulation for an oriented compact surface arises as the contact graph of a circle packing on a geometrized surface (\cite[Section~13.7]{Thurston}; see also~\cite{ColindeVerdiere}). In the case of a closed Riemann surface of genus greater than one, one obtains a hyperbolic metric which realizes given combinatorial structure by a circle pattern on the surface. The hyperbolic metric obtained is unique up to isometries, and this theorem is also called \emph{discrete uniformization theorem}~\cite[Chapter~4]{Stephenson}. In this article, we offer yet another way of endowing a hyperbolic metric on a surface in a manner that given weighted graph is realized as the image of a harmonic map with the least energy among a fixed homotopy class as well as all the hyperbolic metrics.



Let $X=(V, E, m_E)$ be a finite weighted graph with $m_E: E \to (0, \infty)$. We identify each edge $e$ with the unit interval $[0, 1]$, and denote by $\wbar e$ its reversed edge. We understand $E$ the set of all oriented edges $e$, $\wbar e \in E$, and the weight function $m_E$ is symmetric, i.e., $m_E(e)=m_E(\wbar e)$.


For any closed Riemann surface $(S, G)$ equipped with a hyperbolic metric $G$, let $f:X \to (S, G)$ be a piecewise smooth map, i.e., the restriction $f_e:[0, 1] \to (S, G)$ is piecewise smooth for each edge $e$. We define the \emph{energy} of $f$ by
\begin{equation}\label{Eq:Dirichlet}
E_G(f):=\frac{1}{2}\sum_{e\,\in\,E}m_E(e)\int_0^1\left\|\frac{d f_e}{dt}(t)\right\|_G^2\,dt.
\end{equation}
It is known that the energy $E_G(f)$ attains the minimum by a harmonic map $h_G$ among the set of all piecewise smooth maps homotopic to $f$ (Section~\ref{Sec:harmonic}). A \emph{harmonic map} $h$ from $X$ into a hyperbolic surface $(S, G)$ is a map such that $h_e$ is a (constant speed) geodesic for each $e \in E$, and satisfies the \emph{balanced condition}
\[
\sum_{e\,\in\,E_x}m_E(e)\frac{d h_e}{dt}(0)=0, \quad \text{for each }x \in V,
\]
where $E_x=\{e \in E \ : \ e_0=x\}$, the set of edges emanating from $x$. Moreover, if a harmonic map $h_G$ is not homotopic to a point nor to a closed curve in $(S, G)$, then $h_G$ is unique as a map due to the negative curvature of $G$. Then, regarding the energy $E_G(h_G)$ of the harmonic map $h_G$ as a function of hyperbolic metrics $G$ on $S$, we may further minimize $E_G(h_G)$ given a fixed homotopy class $\Cc$ of $f$. A central question we consider is whether the energy $E_G(f)$ attains the minimum for a pair $(h_0, G_0)$ of a map $h_0: X \to S$ in $\Cc$ and a hyperbolic metric $G_0$, or not. We show that under some necessary topological condition on $f:X \to S$, the joint minimum of $E_G(f)$ does exist; namely, we solve the double minimization problem where the first minimization is the classical calculus of variation and the second minimization is the variation among the hyperbolic metrics. Furthermore, we identify the hyperbolic metric $G$ when the underlying weighted graph $X$ admits an automorphism group which is compatible with the mapping class group of $S$.

Note that for $i=1, 2$, any two pairs $(h_i, G_i)$ of a hyperbolic metric $G_i$ on $S$ and a harmonic map $h_i:X \to (S, G_i)$ in $\Cc$ have the same energy $E_{G_1}(h_{1})=E_{G_2}(h_{2})$ if there exists an isometry $\f: (S, G_1) \to (S, G_2)$ homotopic to the identity map on $S$ and $h_2=\f \circ h_1$. It holds that if a pair $(h, G)$ attains the minimum of the energy $E_G(h)$, then the hyperbolic metric $G$ on $S$ is unique up to isometries homotopic to the identity map on $S$ under some condition on the fixed homotopy class $\Cc$ as we see below.

We say that a continuous map $f:X \to S$ \emph{fills} $S$ if it is injective and the complement of the image $f$ is a disjoint union of (topological) disks. For example, any skeleton of triangulation of $S$ gives rise to such $X$ and $f$. Let us denote the set of all piecewise smooth maps homotopic to $f$ by $\Cc=[f]$, and the space of hyperbolic metrics on $S$ by $\Mc_{-1}(S)$. We define
\[
\Ec: \Cc \times \Mc_{-1}(S) \to [0, \infty), \qquad (f, G) \mapsto E_G(f).
\]
Note that for any piecewise smooth map $f: X \to (S, G)$, we have $E_G(f)<\infty$, and any continuous map $f$ is homotopic to a piecewise smooth map.

\begin{theo}\label{Thm:main}
Let $X=(V, E, m_E)$ be a connected finite weighted graph, and $S$ be a closed Riemann surface of genus greater than one. If $f:X \to S$ is a continuous map which fills $S$, then there exists a pair $(h_0, G_0)$ which attains the minimum of $\Ec$ on $\Cc \times \Mc_{-1}(S)$, where $\Cc$ is the homotopy class of $f$. Moreover, the hyperbolic metric $G_0$ on $S$ is unique up to isometries homotopic to the identity map on $S$.
\end{theo}



In fact, we prove a more general theorem when $f$ induces a surjective homomorphism $f_\ast: \pi_1(X, x_0) \to \pi_1(S, f(x_0))$ in Theorem~\ref{Thm:general_main}. The unique hyperbolic metric $G_0$ in Theorem~\ref{Thm:main} realizes the least energy harmonic embedding of $X$ into the hyperbolic surface $(S, G_0)$ determined by the weighted graph $X$ and the fixed homotopy class of maps $f: X \to S$.

The problem of minimizing energy~\eqref{Eq:Dirichlet} was studied by Colin de Verdière~\cite{ColindeVerdiereComment} for the embedding of finite weighted graphs into surfaces. The problem of minimizing energy also in hyperbolic metrics was suggested by Kotani and Sunada~\cite[p.~7, Section~2]{KotaniSunadaStandard} and has been reiterated by Sunada~\cite[p.~124, Section~7.7]{SunadaTopological}. They showed the corresponding result in the case of flat tori (of dimension at least $2$); for a connected finite weighted graph $X$, and for an $n$-dimensional torus $\T^n$, if $f:X \to \T^n$ induces a surjective homomorphism from $\pi_1(X)$ to $\pi_1(\T^n)$, then there exists a pair of flat metric $G_0$ on $\T^n$ and a harmonic map $h_0: X \to (\T^n, G_0)$ homotopic to $f$ such that the corresponding energy functional $\Ec$ attains the minimum on $\Cc \times \Mc_0(\T^n)$, where $\Cc=[f]$ and $\Mc_0(\T^n)$ is the space of flat metrics on $\T^n$ whose volume is normalized to $1$ (and such a pair is unique up to translations on $\T^n$). In fact, they proved the result by constructing a flat metric explicitly by using the $1$-homology group of graph $X$, and called the resulting harmonic map $h_0$ into $(\T^n, G_0)$ (or, its equivariant lift on an abelian covering of $X$ into $\R^n$) the \emph{standard realization} of $X$. Then, they proposed the problem in the case of closed hyperbolic surfaces whether such an energy functional has a minimum as a non-Euclidean analogue to their result in dimension~$2$. It seems that a direct construction of such a pair of a metric and harmonic map into closed hyperbolic surfaces would not work by a non-linear nature of target spaces.


Our approach to show the existence of a minimum pair of energy functional $\Ec$ for surfaces is that we define a function $\Ec_\Cc$ associated to $\Ec$ on the Teichm\"uller space $\Tc(S)$. Namely, since the energy $E_G(f)$ attains the minimum at a harmonic map $h_G$ in the homotopy class $\Cc$ of $f$, the function $G \mapsto E_G(h_G)$ is defined on the space of hyperbolic metrics $\Mc_{-1}(S)$. This function naturally yields a function on the Teichm\"uller space $\Tc(S)$ of $S$. Recalling that the Teichm\"uller space $\Tc(S)$ consists of isotopy classes of pairs $((\SS, G_{\SS}), \f)$ of a hyperbolic surface $(\SS, G_{\SS})$ and a diffeomorphism (a marking) $\f:S \to \SS$, we write $G=\f^\ast G_\SS$ the pull-back metric of $G_\SS$ on $S$. Then, the function
\[
\Ec_\Cc([G]):=E_{G}(h_G),
\]
is well-defined on $\Tc(S)$ (Section~\ref{Sec:energy}). We show that the function $\Ec_\Cc$ on $\Tc(S)$ is strictly convex with respect to the Weil--Petersson metric (Theorem~\ref{Thm:convex}) and proper if $f$ induces a surjective homomorphism from $\pi_1(X)$ to $\pi_1(S)$ (Theorem~\ref{Thm:proper}). This shows that a minimum of the function $\Ec_\Cc$ on $\Tc(S)$ exists uniquely and we show that it also yields a minimum pair of the original energy functional $\Ec$.



\section{The energy functional on the Teichm\"uller space}\label{Sec:energyfunctionalTc}

\subsection{Discrete harmonic maps}\label{Sec:harmonic}

Let $S$ be a closed Riemann surface of genus greater than one endowed with a hyperbolic metric $G$, and $X=(V, E, m_E)$ be a finite weighted graph. Theorems by Colin de Verdière~\cite{ColindeVerdiereComment} and by Kotani--Sunada in~\cite{KotaniSunadaStandard} imply that for any continuous map $f:X \to S$, the minimum of energy $E_G$ is achieved by a harmonic map $h$ among all piecewise smooth maps $\Cc$ homotopic to $f$. Moreover, $h$ is a harmonic map in $\Cc$ if and only if $h$ attains the minimum of the energy $E_G$ in $\Cc$. We give a proof adapted to our setting for the sake of convenience.



\begin{theo}[{\cite[Théorème~1]{ColindeVerdiereComment} and~\cite[Theorems~2.2, 2.3 and~2.5]{KotaniSunadaStandard}}]\label{Thm:KSharmonic}
Fix a hyperbolic metric $G$ in $S$ and a (not necessarily connected) finite weighted graph $X=(V, E, m_E)$. For any continuous map $f:X \to (S, G)$, let $\Cc=[f]$ be the set of all piecewise smooth maps homotopic to $f$.
\begin{enumerate}\romanenumi
\item \label{theo2.1.1} There exists a map $h$ such that the energy $E_G$ is the minimum in $\Cc$. Moreover, if a map $h$ attains the minimum of $E_G$ in $\Cc$, then $h$ is a harmonic map.
\item \label{theo2.1.2}For any harmonic maps $h_0$ and $h_1$ in $\Cc$, we have $E_G(h_0)=E_G(h_1)$.
\item \label{theo2.1.3}If the image of $f$ on each connected component of $X$ is not homotopic to a point nor a closed circle, then there exists a unique harmonic map $h$ in $\Cc$ (as a map). Moreover, the Hessian ${\Hess}_{E_G}$ of the energy functional $E_G$ at $h$ is non-degenerate.
\end{enumerate}
\end{theo}

%\proof
\begin{proof}
Let $\Cc_\geo$ be the set of piecewise geodesic maps in $\Cc$, i.e., maps restricted to each edge $e$ is a geodesic. Note that $\Cc_\geo$ is non-empty because we find a map $f_\geo \in \Cc_\geo$ for any piecewise smooth map $f\in \Cc$ so that $f(x)=f_\geo(x)$ for any $x\in V$, and furthermore $E_G(f_\geo)\leq E_G(f)$. If we endow the $C^1$-topology on $\Cc_\geo$, then the energy functional $E_G$ restricted to $\Cc_\geo$ is continuous and proper (i.e., for each $R \ge 0$, the set $\{f \in \Cc_\geo \ : \ E_G(f) \le R\}$ is compact) on $\Cc_\geo$. Thus, there exists a piecewise geodesic map $h$ which attains the minimum of $E_G$ in $\Cc$.


Note that a map $h:X \to (S, G)$ is a harmonic map if and only if for any smooth variations $f_s$ for $s\in (-\e, \e)$, $\e>0$ and $f_0=h$, we have $dE_G(f_s)/ds|_{s=0}=0$ by the first variation formula~\eqref{Eq:firstVF} in Lemma~\ref{Lem:first and second} since $h$ satisfies that
\[
\nabla_{\p_t h_{e}}\p_t h_{e}=0 \quad \text{for any }e \in E \quad \text{and} \quad \sum_{e\,\in\,E_x}m_E(e)\p_t h_{e}(0)=0 \quad \text{for any }x \in V,
\]
where we write $\p_t h_{e}(t):=(d h_e/dt)(t)$ for $e \in E$ and $t \in [0, 1]$. Hence if $h$ attains the minimum of $E_G$ in $\Cc$, then such $h$ is always a harmonic map. This proves~\eqref{theo2.1.1}.


For any two harmonic maps $h_0$ and $h_1$ in $\Cc$, there exists a homotopy $f_s: X \to S$ so that $f_s(x)$ for $s \in [0, 1]$ is a (unique) geodesic from $h_0(x)$ to $h_1(x)$ in the hyperbolic surface $(S, G)$ for any $x\in X$ by perturbing the given homotopy if necessary. We denote the variational vector field of $f_s$ by $V_e^s(t):=\tfrac{d}{ds} f_{s,\,e}(t)$ and $T_e^s:=\p_t f_{s,\,e}(t)$. Then, $\nabla_{V_e^s}V_e^s=0$ for the Levi-Civita connection $\nabla$, and hence, the second variation formula~\eqref{Eq:secondVF} in Lemma~\ref{Lem:first and second} shows that
\begin{equation}\label{positive}
\frac{d^2}{ds^2}E_G(f_s)=\sum_{e\,\in\,E} m_E(e)\int_0^1\left\{\left\|\nabla_{T_e^s} V_e^s\right\|_G^2+\left\|\left(V_e^s\right)^\perp\right\|_G^2\left\|T_e^s\right\|_G^2\right\}\,dt\geq 0
\end{equation}
for any $s\in [0,1]$ since $(S,G)$ has constant sectional curvature $-1$, where $(V_e^s)^\perp$ means the orthogonal projection of $V_e^s$ normal to $T_e^s$ if $T_s^s \neq 0$, and $0$ otherwise. Since $h_0$ and $h_1$ are harmonic maps, we have $dE_G(f_s)/ds|_{s=0}=dE_G(f_s)/ds|_{s=1}=0$, and~\eqref{positive} implies that $dE_G(f_s)/ds\equiv 0$ on $[0,1]$ and $E_G(h_0)=E_G(h_1)$. This shows~\eqref{theo2.1.2}.

Furthermore, in fact, we have $d^2E_G(f_s)/ds^2\equiv 0$ on $[0, 1]$, and by~\eqref{positive}, $\nabla_{T_e^s}V_e^s=0$ and since $(S, G)$ has negative sectional curvature, $V_e^s$ and $T_e^s$ are proportional for any $s \in [0, 1]$. If the image of $f$ on each connected component $X^0$ of $X$ is not homotopic to a point nor a closed circle, then for any $s \in [0, 1]$ and in each connected component $X^0=(V^0, E^0)$, there exists a vertex $x \in V^0$ for which $\{T_e^s\}_{e\,\in\,E^0_x}$ contains at least two linearly independent vectors in $T_{f_s(x)}S$ since $f_s$ is homotopic to $f$; and thus $V_e^s\equiv 0$ for all $e \in E^0_x$. In fact, $V_e^s \equiv 0$ for all $e \in E$ by parallel transports along the image of edges since $\nabla_{T_e^s}V_e^s=0$ and each $X^0$ is connected, and this holds for any $s \in [0, 1]$. Therefore $h_0=h_1$. Moreover, ${\Hess}_{E_G}(V, V)=(d^2/ds^2)|_{s=0}E_G(f_s)=0$ implies $V_e \equiv 0$ for all $e \in E$. We conclude the proof of~\eqref{theo2.1.3}.
\end{proof}



The following proposition guarantees that for a given closed hyperbolic surface $(S, G)$ and a harmonic map $h:X \to (S, G)$, harmonic maps $h_s$ depends smoothly on a smooth change of hyperbolic metrics $G_s$. We give the proof in a more general setting in Proposition~\ref{smooth dependence} in Appendix~\ref{AppB}.

\begin{prop}\label{Prop:smooth}
Let $X$ be a connected finite weighted graph, $(S, G)$ be a closed hyperbolic surface and $h: X \to (S, G)$ be a harmonic map. Suppose that the image of $h$ is not a point nor a closed geodesic. Then, for any $\e>0$ and for any smooth family of hyperbolic metrics $\{G_s\}_{s\,\in\,(-\epsilon,\,\epsilon)}$ with $G_0=G$ in $\Mc_{-1}(S)$, there exists a unique family of maps $h_s:X \to S$ for $s\in (-\epsilon,\epsilon)$ with $h_0=h$ such that $h_s:X \to (S, G_s)$ is a harmonic map for every $s \in (-\e, \e)$ and $h_s$ is smooth with respect to the variable $s$.
\end{prop}

\begin{proof}
If the image of $h$ is not homotopic to a point nor a closed geodesic, then Theorem~\ref{Thm:KSharmonic}$\MK$\eqref{theo2.1.3} implies that for each $s \in (-\e, \e)$ there exists a unique harmonic map $h_s: X \to (S, G_s)$ homotopic to $h$ and the Hessian of $E_G$ at $h_s$ is non-degenerate. Hence by Proposition~\ref{smooth dependence}, since $G_s \in \Mc_{-1}(S)$ for all $s \in (-\e, \e)$, for each $G_s$ and for the harmonic map $h_s: X \to (S, G_s)$, there exists an open set $\Uc$ around $G_s$ in $\Mc_{-1}(G)$ and a smooth map $\wt h: \Uc \to C^\infty(X, S)$ (the space of piecewise smooth map from $X$ into $S$) such that $\wt h(G'):X \to (S, G')$ is a harmonic map for all $G' \in \Uc$ and $\wt h(G_s)=h_s$. Since for each $s \in (-\e, \e)$, the harmonic map $h_s:X \to (S, G_s)$ is unique in the homotopy class, covering the curve $\{G_s\}_{s\,\in\,(-\e,\,\e)}$ by such open sets, we obtain the claim.
\end{proof}


\subsection{The energy functional on the Teichm\"uller space}\label{Sec:energy}
Recall that the Teichm\"uller space of $S$ consists of equivalence classes of pairs $((\SS, G_{\SS}), \f)$ of a hyperbolic surface $(\SS, G_{\SS})$ and a diffeomorphism (a marking) $\f:S \to \SS$, where $((\SS_1, G_1), \f_1) \sim ((\SS_2, G_2), \f_2)$ if and only if $\f_2 \circ \f_1^{-1}$ is homotopic to an isometry from $(\SS_1, G_1)$ to $(\SS_2, G_2)$. We discuss the standard topology in $\Tc(S)$ (see e.g., \cite[Section~10.3, p.~269]{FarbMargalit}). The mapping class group $\Mod(S)$ is the group of isotopy classes of orientation-preserving diffeomorphisms on $S$. The group $\Mod(S)$ acts on the Teichm\"uller space by the change of markings
\[
[\phi]\cdot [(\SS, G_{\SS}), \f]:=\left[(\SS, G_{\SS}), \f \circ \phi^{-1}\right] \qquad \text{for }\;[\phi] \in \Mod(S),
\]
and the action is properly discontinuous~\cite[Theorem~12.2, p.~350]{FarbMargalit}. The moduli space $\Mc(S)$ of $S$ is the quotient space of $\Tc(S)$ by the action of the mapping class group $\Mod(S)$.



The energy $E_G(f)$ attains the minimum at a harmonic map $h_G$ in the homotopy class $\Cc$ of $f$. The function $G \mapsto E_G(h_G)$ on the space of hyperbolic metrics $\Mc_{-1}(S)$ naturally defines a function on the Teichm\"uller space $\Tc(S)$ of $S$. For a point $[(\SS, G_{\SS}), \f]$ of the Teichm\"uller space $\Tc(S)$, we write $\f^\ast G_\SS$ the pull-back metric of $G_\SS$ on $S$ by $\f:S \to (\SS, G_\SS)$. If $h_1: X \to (S, \f_1^\ast G_1)$ is harmonic and $((\SS_1, G_1), \f_1) \sim ((\SS_2, G_2), \f_2)$, then the identity map on $S$ is homotopic to an isometry $\i: (S, \f_1^\ast G_1) \to (S, \f_2^\ast G_2)$, and $\i \circ h_1:X \to (S, \f_2^\ast G_2)$ is also a harmonic map homotopic to $h_1$; hence for any harmonic map $h_2: X \to (S, \f_2^\ast G_2)$ homotopic to $h_1$, one has
\[
E_{\f_1^\ast G_1}(h_1)=E_{\f_2^\ast G_2}(h_2)=\min_{h\,\in\,\Cc}E_G(h),
\]
by Theorem~\ref{Thm:KSharmonic}$\MK$\eqref{theo2.1.1} and~\eqref{theo2.1.2}. We define
\[
\Ec_\Cc([G]):=E_{G}(h_G),
\]
where $G=\f^\ast G_{\SS}$ is the hyperbolic metric induced on $S$ by a pair $((\SS, G_{\SS}), \f)$ representing a point in $\Tc(S)$ and $h_G$ is a harmonic map from $X$ to $(S, G)$ in the homotopy class $\Cc=[f]$. Note that $\Ec_\Cc$ is a function on the Teichm\"uller space
\[
\Ec_\Cc:\Tc(S) \to [0, \infty), \quad \text{and} \quad \Ec_\Cc([G])=E_{G}(h_G)=\min_{h\,\in\,\Cc}E_G(h).
\]
We show that the function $\Ec_\Cc$ on $\Tc(S)$ is strictly convex with respect to the Weil--Petersson metric; we will give the proof in Section~\ref{Sec:convex}.

\begin{theo}\label{Thm:convex}
Let $X=(V, E, m_E)$ be a connected finite weighted graph and $S$ be a closed Riemann surface of genus greater than one. For any homotopy class $\Cc$ of piecewise smooth maps $f: X \to S$ such that the image of $f$ is not homotopic to a point nor a closed geodesic, the function $\Ec_\Cc([G])=\min_{f \,\in\,\Cc}E_G(f)$ on $\Tc(S)$ is strictly convex with respect to the Weil--Petersson metric.
\end{theo}



Recall that a continuous map $f:X \to S$ fills $S$ if there exists an injective map $f_1:X \to S$ homotopic to $f$ such that the complement of the image $f_1$ is a disjoint union of disks.

\begin{lemm}\label{Lem:fill}
If $f:X \to S$ fills $S$, then the induced homomorphism $f_\ast: \pi_1(X, x_0) \to \pi_1(S, f(x_0))$ is surjective.
\end{lemm}

%\proof
\begin{proof}
There exists an injective map $f_1$ homotopic to $f$ such that the complement of the image $f_1$ is a disjoint union of disks. Then, for each disk, choosing a point from the interior, we take a homotopy on the disk such that it pushes outside of a small neighborhood of the point to the boundary and remains identity on the boundary. Patching these homotopies together, we obtain a homotopy on the surface.

For any point $x_0$ in $X$, and any loop $\g$ based at $f_1(x_0)$ in $S$, up to a small perturbation of $\g$ if necessary, composing the homotopy constructed, we obtain a loop confined in the boundaries of disks. Since $f_1$ is injective, there is a loop $c$ in the graph $X$ based at $x_0$, whose image by $f_1$ is homotopic to the original loop $\g$ relative to $f_1(x_0)$ on the surface. This shows that the loop $c$ satisfies $f_{1\,\ast} [c]=[\g]$. Thus, $f_{1 \ast}: \pi_1(X, x_0) \to \pi_1(S, f_1(x_0))$ is surjective, and since $f_1$ is homotopic to $f$, we conclude the claim.
\end{proof}


If we fix the homotopy class $\Cc$ of $f$ such that $f$ induces a surjective homomorphism from $\pi_1(X)$ to $\pi_1(S)$, then the energy functional $\Ec_\Cc$ is proper; we shall give the proof in Section~\ref{Sec:proper}.

\begin{theo}\label{Thm:proper}
Let $X=(V, E, m_E)$ be a finite weighted graph, and $S$ be a closed Riemann surface of genus greater than one. Let $\Cc$ be the homotopy class of a continuous map $f:X \to S$ such that the induced homomorphism $f_\ast:\pi_1(X, x_0) \to \pi_1(S, f(x_0))$ is surjective. Then, the function $\Ec_\Cc:\Tc(S) \to [0, \infty)$ is proper, i.e., for any $R \ge 0$, the sublevel set of $\Ec_\Cc$,
\[
\big\{\left[(\SS, G_\SS), \f\right] \in \Tc(S) \, : \, \Ec_\Cc([G]) \le R, \ G=\f^\ast G_\SS \big\}
\]
is compact relative to the standard topology in $\Tc(S)$.
\end{theo}

In general, we are not able to remove the condition that $f$ induces a surjective homomorphism from $\pi_1(X)$ to $\pi_1(S)$ in Theorem~\ref{Thm:proper}. For example, if we have a simple closed curve in $S$ as an image of $f$, then one is able to make the length of closed geodesic homotopic to $f$ arbitrary small by changing hyperbolic metrics in $S$.



\begin{theo}\label{Thm:general_main}
Let $X=(V, E, m_E)$ be a finite weighted graph, and $S$ be a closed Riemann surface of genus greater than one. If $\Cc=[f]$ is the homotopy class of $f$ which induces a surjective homomorphism from $\pi_1(X)$ to $\pi_1(S)$, then the energy functional $\Ec_\Cc$ has a unique minimum point in the Teichm\"uller space $\Tc(S)$.
\end{theo}


\begin{proof}
Theorem~\ref{Thm:proper} implies that there exists a minimum of $\Ec_\Cc$ on $\Tc(S)$ if $f$ induces a surjective homomorphism from $\pi_1(X)$ to $\pi_1(S)$. Since any two points in $\Tc(S)$ can be connected by a Weil--Peterson geodesic by a theorem by Wolpert~\cite[Corollary~5.6]{WolpertNielsen}, Theorem~\ref{Thm:convex} shows that a minimum is unique.
\end{proof}


\begin{proof}[{Proof of Theorem~\ref{Thm:main}}]
Lemma~\ref{Lem:fill} and Theorem~\ref{Thm:general_main} implies that there exists a unique point $[(\SS, G_{\SS}), \f]$ in the Teichm\"uller space $\Tc(S)$ such that $\Ec_\Cc$ attains the minimum. Since $\Ec_\Cc([G])=E_{G}(h_G)=\min_{h\,\in\,\Cc}E_G(h)$ for a harmonic map $h_G:X \to (S, G)$ in the homotopy class $\Cc$ and $G=\f^\ast G_{\SS}$, the pair $(h_G, G)$ of a harmonic map $h_G$ in $\Cc$ and the hyperbolic metric $G$ on $S$ attains the minimum $E_G(h_G)$. Note that the hyperbolic metric $G \in \Mc_{-1}(S)$ is unique up to isometries homotopic to the identity on $S$, we conclude the claim.
\end{proof}



\section{Convexity of the energy functional}\label{Sec:convex}

In this section, we compute the second derivative of the energy functional $\Ec_\Cc$ for $\Cc=[f]$. The argument is inspired by Wolf, who showed the convexity of the \emph{length} functional of a \emph{closed} curve along the Weil--Petersson geodesic~\cite{Wolf}. We discuss the energy of curves and extend it to any finite weighted graphs. First, we recall some differential geometric aspects and facts that we use on the Teichm\"uller space, following~\cite{Yamada}, \cite{Yamada2} and~\cite{Wolf}.


\subsection{Hessian formula of the energy functional}\label{ss1}


Let $\Met(S)$ be the set of all smooth metrics on $S$. Then, $\Mc_{-1}(S)$ is the subset of $\Met(S)$ consisting of metrics of constant sectional curvature $K=-1$. The tangent spaces of $\Met(S)$ is regarded as
\[
T_G\Met(S)=\left\{\text{symmetric}\ (0,2)\text{-tensors on}\ S\right\}.
\]
The {Teichm\"uller space} $\Tc(S)$ of $S$ is identified with the quotient space of $\Mc_{-1}(S)$ by ${\Diff}_0(S)$ via the map
\[
\Tc(S) \to \Mc_{-1}(S)/{\Diff}_0(S), \qquad \left[\left(\Sigma, G_\SS\right), \varphi\right]\mapsto \left[\varphi^\ast G_\SS\right],
\]
where ${\Diff}_0(S)$ is the identity component of the group of diffeomorphisms on $S$.



The deformations arising from ${\Diff}_0(S)$ are given by the set of Lie derivatives $L_XG$ of $G$ for smooth vector fields $X$ on $S$. The tangent space to $\Tc(S)$ at $[G]$ is identified with
\[
T_{[G]}\Tc(S)=\big\{Q\in T_G\Met(S);\ \tr_GQ=0\quad \text{and}\quad \delta_GQ=0\big\},
\]
where $\delta_GQ:=-{\tr}_{12}(\nabla Q)=-G^{ij}\nabla{Q}(\p_i,\p_j, \cdot)$ for the Levi--Civita connection $\nabla$ of $(S, G)$~\cite[Section~3.2]{Yamada2}, where for a local coordinate $(x^1,x^2)$ on $S$, we abbreviate $\p/\p x^1$ and $\p/\p x^2$ by $\p_1$ and $\p_2$, respectively. We denote the matrix representation of any symmetric $(0,2)$-tensor $Q$ by $Q=(Q_{ij})_{i,\,j\,=\,1,\,2}=(Q(\p_i,\p_j))_{i,\,j\,=\,1,\,2}$.

\begin{lemm}\label{Lem:Qholomorphic}
Let $(S,G)$ be a hyperbolic surface of constant sectional curvature $K=-1$ and $Q$ be a symmetric $(0,2)$-tensor on $S$ such that ${\tr}_G Q=0$ and $\delta_GQ=0$. Then, we have the following:
\begin{enumerate}\romanenumi
\item\label{lemm3.1.1} For any isothermal chart $(U, z=x^1+\sqrt{-1}x^2)$ of $(S,G)$, the function $Q_{11}-\sqrt{-1}Q_{12}$ is holomorphic on $U$.
\item\label{lemm3.1.2} It holds that
\[
\left(\Delta_G+4\right)\|Q\|^2_G\geq 0,
\]
where $\Delta_G$ is the Laplace--Beltrami operator acting on $C^{\infty}(S)$.
\end{enumerate}
\end{lemm}

 
\begin{proof}
Take an isothermal coordinate $(U, z=x^1+\sqrt{-1}x^2)$ so that the metric $G$ is expressed by $G(z)=\rho(z)|dz|^2=\rho(x^1, x^2)((dx^1)^2+(dx^2)^2)$ for a positive real function $\rho$ on $U$. Note that the Christoffel symbol $\Gamma_{ij}^k$ for the Levi--Civita connection on any isothermal coordinate satisfies
\begin{align*}
\begin{cases}
\Gamma_{11}^1=\Gamma_{12}^2=\Gamma_{21}^2=-\Gamma_{22}^1=\frac{1}{2}\p_1\log {\rho},\\
\Gamma_{22}^2=\Gamma_{21}^1=\Gamma_{12}^1=-\Gamma_{11}^2=\frac{1}{2}\p_2\log{\rho},\\
\end{cases}
\end{align*}
Then we have
\begin{align*}
(\delta_{G}Q)(\p_k)=-\rho^{-1}(Q_{1k,\,1}+Q_{2k,\,2}),
\end{align*}
for $k=1,2$ on $U$, where $Q_{ij,\,k}:=\p_kQ_{ij}$. Using $\delta_{G}Q=0$, $Q_{22}=-Q_{11}$ and $Q_{12}=Q_{21}$, we obtain
\[
Q_{11,\,1}=-Q_{12,\,2}\quad \text{and}\quad Q_{11,\,2}=Q_{12,\,1},
\]
hence, $Q_{11}-\sqrt{-1}Q_{12}$ is a holomorphic function on $U$ since it satisfies the Cauchy--Riemann equation. This shows~\eqref{lemm3.1.1}.


Let us show~\eqref{lemm3.1.2}. If $\|Q\|_G^2\equiv 0$ or equivalently $Q\equiv 0$, then the statement is obvious. We may assume the zeros of the function $\|Q\|_G^2$ is isolated since $Q_{11}-\sqrt{-1}Q_{12}$ is holomorphic by (i). Take any $p\in S$ with $\|Q\|_G^2(p)\neq 0$ and an isothermal coordinate $(U, z=x^1+\sqrt{-1} x^2)$ around $p$ so that $G=\rho((dx^1)^2+(dx^2)^2)$ and $\|Q\|_G^2\neq 0$ on $U$. Note that, we have $\Delta_G=\rho^{-1}(\p_1^2+\p_2^2)$ in the isothermal coordinate.

We have
\[
\|Q\|_G^2=2\rho^{-2}\left(Q_{11}^2+Q_{12}^2\right)
\]
since $Q$ is symmetric and trace free tensor. Taking the logarithm of this equation, we have $\log\|Q\|_G^2=\log|Q|^2-2\log\rho+\log2$, where we set $|Q|^2:=|Q_{11}-\sqrt{-1}Q_{12}|^2=Q_{11}^2+Q_{12}^2$, and hence, we obtain
\begin{align}\label{dlog}
\Delta_G\log\|Q\|_G^2&=\Delta_{G}\log|Q|^2-2\Delta_G\log\rho
\end{align}
Since $Q_{11}-\sqrt{-1}Q_{12}$ is holomorphic on $U$ by~\eqref{lemm3.1.1}, we see $\Delta_{G}\log|Q|^2=0$ on $U$. On the other hand, in the isothermal coordinate, the sectional curvature $K$ is given by $K=-(1/2)\Delta_G\log \rho$. Since we assume $K=-1$, the equation~\eqref{dlog} becomes
\[
\Delta_G\log\|Q\|_G^2=-4.
\]
By using the formula $\Delta_G\log F=\tfrac{\Delta_G F}{F}-\tfrac{|\nabla F|_G^2}{F^2}$, we obtain
\[
\frac{\Delta_G\|Q\|_G^2}{\|Q\|_G^2}+4=\frac{|\nabla \|Q\|_G^2|_G^2}{\|Q\|_G^4}\geq 0,
\]
as required, and conclude the proof of~\eqref{lemm3.1.2}.
\end{proof}


\begin{rema}\label{Rem:QD}
Given an isothermal coordinate $z=x^1+\sqrt{-1}x^2$ on $(S,G)$ with $G(z)=\rho(z)|dz|^2$, we define $\Phi:=(Q_{11}-\sqrt{-1}Q_{12})dz^2$. It turns out that $\Phi$ defines a global section of $(T^{1,\,0}S)^*\otimes (T^{1,\,0}S)^*$ and the section is holomorphic by Lemma~\ref{Lem:Qholomorphic}$\MK$\eqref{lemm3.1.1}. The section $\Phi$ is a \emph{holomorphic quadratic differential} on the Riemann surface $S$. One verifies that there exists a one-to-one correspondence between $Q$ and $\Phi$.
\end{rema}



We define an $L^2$-pairing on $T_G\Met(S)$ by
\[
\left\langle Q_1, Q_2\right\rangle:=\int_{S} \left\langle Q_1(x), Q_2(x)\right\rangle_{G(x)}\,d\mu_G(x) \quad \text{for }Q_1,Q_2\in T_G\Met(S),
\]
where $\m_G$ is the area form with respect to $G$. It induces an $L^2$-pairing on $T_{[G]}\Tc(S)$ and a Riemannian metric on $\Tc(S)$. This metric is called the \emph{Weil--Petersson metric} on $\Tc(S)$. Then, the mapping class group $\Mod(S)$ acts on $\Tc(S)$ as isometries of the Weil--Petersson metric.



Let $[G_{s}]\in \Tc(S)$ be a geodesic relative to the Weil--Petersson metric for $s \in (-\epsilon, \epsilon)$ and for $\e>0$. Then, we can lift $\{[G_{s}]\}_{s\,\in\,(-\e,\,\e)}$ to a horizontal geodesic $\{G_{s}\}_{s\, \in\,(-\e,\,\e)}$ in $\Mc_{-1}(S)$ since $\Mc_{-1}(S)\to \Tc(S)$ is a Riemannian submersion (\cite[Proposition~2.3.1]{Yamada}, and see also~\cite[Section~3.6]{Yamada2}). By Proposition~\ref{Prop:smooth}, there exists a family of harmonic maps $h_{s}: X \to (S, G_s)$ for $s \in (-\e, \e)$ associated with $G_{s}$ in the fixed homotopy class $\Cc=[f]$ such that $h_{s}$ is differentiable with respect to $s \in (-\e, \e)$. Let us define the associated variational vector field
\[
V_e(s,t):=\frac{\p}{\p s}h_{s,\,e}(t) \quad \text{for each edge }e\in E,
\]
where $h_{s,\,e}:[0, 1] \to S$ is the restriction of $h_s$ on the edge $e$.


Recall that the energy functional $\Ec_{\Cc}$ for $[G_{s}]$ is given by
\[
\Ec_{\Cc}\left(\left[G_{s}\right]\right):=E_{G_{s}}(h_{s})=
\frac{1}{2}\sum_{e\,\in\,E}m_E(e)\int_{0}^1\left(h_{s}^*G_{s}\right)\left(\p_{t},\p_{t}\right)dt,
\]
where and henceforth we use the short hand notation $\p_t=d/dt$ on $[0, 1]$ for the sake of simplicity.

We shall consider the first and second derivatives of $\Ec_{\Cc}([G_{s}])$. First, for any fixed $s \in (-\e, \e)$, we have
\begin{align}\label{harm}
0=\frac{\p}{\p u} E_{G_{s}}(h_{u})\Big{|}_{u=s}=\frac{1}{2}\sum_{e\,\in\,E}m_E(e)\int_{0}^1h_{s}^*\left(L_{V_e(s)}G_{s}\right)(\p_{t},\p_{t})dt.
\end{align}
since $h_{s}: X\to (S, G_{s})$ is harmonic for any $s \in (-\e, \e)$ and the energy $E_{G_{s}}(h_{u})$ is minimum at $u=s$ by Theorem~\ref{Thm:KSharmonic}$\MK$\eqref{theo2.1.1} and~\eqref{theo2.1.2}. Therefore we have that
\begin{equation}\label{Eq:harmQ}
\begin{aligned}
\frac{d}{d s}\Ec_{\Cc}([G_{s}])&=\frac{1}{2}\sum_{e\,\in\,E}m_E(e)\int_{0}^1{\p_s}\left(h_{s}^*G_{s}\right)(\p_{t},\p_{t})dt\\
&=\frac{1}{2}\sum_{e\,\in\,E}m_E(e)\int_{0}^1\left\{h_{s}^*\left({\p_s}{G}_{s}\right)+h_{s}^*\left(L_{V_e(s)}G_{s}\right)\right\}\left(\p_{t},\p_{t}\right)dt. \\
&=\frac{1}{2}\sum_{e\,\in\,E}m_E(e)\int_{0}^1h_{s}^*\left({\p_s}{G}_{s}\right)\left(\p_{t},\p_{t}\right)dt,
\end{aligned}
\end{equation}
where we have used~\eqref{harm} in the last equality. The above computation yields a characterization of critical points for $\Ec_\Cc$. We prove the following proposition regarding the first variation of $\Ec_\Cc$ although we will not make use of it in this paper.


\begin{prop}
For any $[G]$ in $\Tc(S)$, $[G]$ is a critical point of $\Ec_\Cc$ on $\Tc(S)$ if and only if any harmonic map $h:X \to (S, G)$ in the homotopy class $\Cc$ satisfies
\[
\sum_{e\,\in\,sE}m_E(e)\int_0^1 \F\left(h_{\ast}\p_t, h_{\ast}\p_t\right)\,dt=0,
\]
for any holomorphic quadratic differentials $\F$ on $(S, G)$.
\end{prop}

\begin{proof}
It follows that $[G] \in \Tc(S)$ is a critical point of $\Ec_\Cc$ if and only if for any harmonic map $h:X \to (S, G)$ in the homotopy class $\Cc$, one has
\[
\sum_{e\,\in\,E}m_E(e)\int_0^1 Q\left(h_{ \ast}\p_t, h_{\ast}\p_t\right)\,dt=0,
\]
for any symmetric $(0, 2)$-tensors $Q$ such that $\tr_{G}Q=0$ and $\d_{G}Q=0$ from~\eqref{harm} and~\eqref{Eq:harmQ}. In any isothermal coordinate $(U, z=x^1+\sqrt{-1}x^2)$ on $(S, G)$, we have
\[
Q\left(h_{\ast}\p_t, h_{\ast}\p_t\right)=\Rel\,\left\{\left(Q_{11}-\sqrt{-1}Q_{12}\right)(h(t))\cdot h'(t)^2\right\},
\]
where we regard $h(t)$ as a complex valued function in the right hand side. Since there is a one-to-one correspondence between $Q \in T_{[G]}\Tc(S)$ and $\F$ holomorphic quadratic differentials on $(S, G)$, where $\F(z)=\f(z)dz^2$ on $U$ for some holomorphic function $\f$ (Remark~\ref{Rem:QD}), and if $\F$ is a holomorphic quadratic differential, then $\sqrt{-1}\F$ is also so; thus we obtain the claim.
\end{proof}



Differentiating~\eqref{Eq:harmQ} in $s$ once more, we obtain
\begin{multline}\label{esvf}
\frac{d^2}{d s^2}\Ec_{\Cc}([G_{s}])\Big{|}_{s=0}\\
=\frac{1}{2}\sum_{e\,\in\,E}m_E(e)\int_{0}^1h_0^*\left\{{\p_s^2}G_{s}\middle|_{s=0}+L_{V_e(0)}\big({\p_s}G_{s}\big|_{s=0}\big)\right\}(\p_{t},\p_{t})dt.
\end{multline}
Here, the second term of the right hand side of~\eqref{esvf} is arranged as follows due to the harmonicity:



\begin{lemm}
\begin{equation}\label{harm3}
\frac{1}{2}\sum_{e\,\in\,E}m_E(e)\int_{0}^1 h_0^*\big(L_{V_e}{\p_s}G_{s}\big|_{s=0}\big)(\p_{t},\p_{t})dt=-\frac{d^2}{d s^2}E_{G_{0}}(h_{s})\Big{|}_{s=0}.
\end{equation}
\end{lemm}



\begin{proof}
Differentiating~\eqref{harm} with respect to the variable $s$ at $s=0$, we have
\begin{equation}\label{e31}
\begin{aligned}
0&=\frac{1}{2}\sum_{e\,\in\,E}m_E(e)\int_{0}^1h^*_0\left\{{\p_s}\left(L_{V_e(s)}G_{s}\right)\middle|_{s=0} +\left(L_{V_e(s)}L_{V_e(s)}G_{s}\right)\middle|_{s=0}\right\}\left(\p_{t},\p_{t}\right)dt\\
&=\frac{1}{2}\sum_{e\,\in\,E}m_E(e)\int_{0}^1h_0^*\left\{\left(L_{V_e}{\p_s}G_{s}\middle|_{s=0}\right)+2\left(L_{V_e}L_{V_e}G_{s}\right)\middle|_{s=0}\right\}\left(\p_{t},\p_{t}\right)dt,
\end{aligned}
\end{equation}
where we used the following identity in the second equality:
\begin{align*}
h_s^*\left({\p_s}L_{V_e(s)}G_{s}\right) &=\frac{\p}{\p s}\frac{\p}{\p u}h_{u+s}^*G_{s}\Big|_{u=0}\\
&=\frac{\p}{\p u}\frac{\p}{\p s}h_{u+s}^*G_{s}\Big|_{u=0}\\
&=\frac{\p}{\p u}\left(h_{u+s}^*\frac{\p}{\p s}G_{s}+h^*_{u+s}L_{V_e}G_{s}\right)\bigg|_{u=0}\\
&=h_{s}^*\left(L_{V_e}\frac{\p}{\p s}G_{s}\right)+h_{s}^*L_{V_e}L_{V_e}G_{s}.
\end{align*}

On the other hand, a similar computation shows that
\begin{multline}\label{e32}
\frac{d^2}{d s^2}E_{G_{0}}(h_{s})\bigg|_{s=0}\\
\begin{aligned}
&=\frac{1}{2}\sum_{e\,\in\,E} m_E(e) \int_0^1 h_0^*\left\{\left(L_{V_e}L_{V_e}G_{0}\right)\middle|_{s=0}+\p_s\left(L_{V_e(s)}G_{0}\right)\middle|_{s=0}\right\}\left(\p_t,\p_t\right)\,dt\\
&=\sum_{e\,\in\,E} m_E(e) \int_0^1 h_0^*\left(L_{V_e}L_{V_e}G_{0}\right)|_{s=0}\left(\p_t,\p_t\right)\,dt,
\end{aligned}
\end{multline}
where we have used $h_0^\ast \p_s(L_{V_e(s)}G_{0})|_{s=0}=h_{0}^*L_{V_e}L_{V_e}G_{0}$ for $h_{0}^*(L_{V_e}\tfrac{\p}{\p s}G_{0})=0$ in the last equality. Combining~\eqref{e31} with~\eqref{e32}, we obtain~\eqref{harm3}.
\end{proof}



Substituting~\eqref{harm3} to~\eqref{esvf}, we have that
\begin{multline}\label{esvf2}
\frac{d^2}{d s^2}\Ec_{\Cc}([G_{s}])\bigg{|}_{s=0} \\
=\frac{1}{2}\sum_{e\,\in\,E}m_E(e)\int_{0}^1h_0^*\left(\frac{d^2}{d s^2}G_{s}\right)\bigg|_{s=0}\left(\p_t,\p_t\right)dt -\frac{d^2}{d s^2}E_{G_{0}}(h_{s})\bigg|_{s=0}.
\end{multline}
The first term in the right hand side of~\eqref{esvf2} can be computed by using the differential geometry of Teichm\"uller space (without any assumption for $h_{s}$).


\begin{prop}[{\cite[Theorem~3.8]{Yamada2} (cf.~\cite[Theorem~2.3.2]{Yamada})}]\label{key1}
For any horizontal lift $\{G_{s}\}_{s\,\in\,(-\e,\,\e)}$ of a geodesic $\{[G_{s}]\}_{s\,\in\,(-\e,\,\e)}$ for $\e>0$ in the Teichm\"uller space $\Tc(S)$ relative to the Weil--Petersson metric with velocity vector $Q:=(\p/\p s)G_s|_{s=0}$, we have
\[
\frac{d^2}{d s^2}G_{s}\bigg|_{s=0}=\left(\frac{1}{4}-\frac{1}{2}\left(\Delta_{G_0}-2\right)^{-1}\right)\|Q\|_{G_0}^2\cdot G_0 +L_{Z}G_0,
\]
for some (auxiliary) vector filed $Z$ on $S$.
\end{prop}


Thus, the rest is to compute the second derivative of ${E}_{G_0}(h_{s})$. For each edge $e\in E$, let
\[
T_e^s(t):=\p_th_{s,\,e}(t),
\]
and $T_e(t)=T_e^0(t)$ when $s=0$ for simplicity.


\begin{defi}\label{Def:V}
For any velocity vector $Q:=(\p_s G_s)|_{s=0}$, we define a $1$-form on $S$ along $h_{0, \,e}(t)$ by
\[
Q_e:=Q\left(T_e, \cdot\,\right) \quad \text{for any }e \in E,
\]
and we denote its $G_0$-metric dual by $Q_e^{\sharp}$, that is, $Q_e^{\sharp}$ satisfies that $G_0(Q_e^{\sharp},\,\cdot)=Q_e$. Furthermore, we define a vector field along $h_{0,\,e}(t)$ by
\begin{equation}\label{VW}
\Vc_e:=
\frac{1}{\|T_e\|_{G_0}}\left(\nabla_{T_e}V_e+\frac{Q_e^{\sharp}}{2}\right)\quad \text{if}\;\|T_e\|_{G_0}\neq 0,
\end{equation}
and $\Vc_e\equiv 0\ \text{if}\; \|T_e\|_{G_0}=0$ for any $e \in E$.
\end{defi}



The following proposition holds for any horizontal lift of a smooth curve in $\Tc(S)$ (not necessarily a Weil--Petersson geodesic).

\begin{prop}\label{key2}
Let $\{[G_s]\}_{s\,\in\,(-\e,\,\e)}$ for $\e>0$ be any smooth curve in $\Tc(S)$. For any horizontal lift $\{G_s\}_{s\,\in\,(-\e,\,\e)}$ of $\{[G_s]\}_{s\,\in\,(-\e,\,\e)}$ with velocity vector $Q:=(\p/\p s)\linebreak G_s|_{s=0}$, and for any smooth family of harmonic maps $h_{s}:X \to (S, G_s)$ for $s \in (-\e, \e)$ with the variational vector field $V$, we have
\begin{align*}
\frac{d^2}{d s^2}E_{G_0}(h_{s})\bigg|_{s=0}&=\frac{1}{2}\sum_{e\,\in\,E}m_E(e)\int_0^1G_0\left(\nabla_{T_e}Q_e^{\sharp}, V_e\right)\, dt\\
&=\frac{1}{2}\sum_{e\,\in\,E}m_E(e)\int_0^1\left(
\frac{1}{4}\|Q\|_{G_0}^2-2\left\|V_e^{\perp}\right\|_{G_0}^2-2\left\|\Vc_e\right\|_{G_0}^2
\right)\cdot \|T_e\|_{G_0}^2\,dt,
\end{align*}
where $V_e^{\perp}$ denotes the $G_0$-normal component of $V_e$ with respect to $T_e$ if $T_e \neq 0$ and $0$ otherwise, and $\Vc_e$ is the vector field on $(S, G_0)$ defined in Definition~\ref{Def:V}.
\end{prop}



\begin{rema}
By the formula given in Lemma~\ref{le3} in the next subsection, the variational vector field $V_e(t)$ along $h_{0,\,e}(t)$ is determined by the velocity vector $Q=(\p_sG_s)|_{s=0}$ and the second order ordinary differential equation
\begin{equation}\label{deqv1}
\nabla_{T_e}\nabla_{T_e}V_e+R\left(V_e, T_e\right)T_e=-\frac{1}{2}\nabla_{T_e}Q_e^{\sharp}
\end{equation}
with initial vector $V_e(x)$ and $\nabla_{T_e}V_e(x)$ for $x=h_{0,\,e}(0)$, where $R$ is the curvature tensor of $(S,G_0)$. Note that the variational vector filed $V$ is determined uniquely if we take a Weil--Petersson geodesic $\{[G_s]\}_{s\,\in\,(-\e, \,\e)}$ and an initial harmonic map $h_0: X\to S$ whose image is not a point nor closed circle by the uniqueness of the harmonic map. Moreover, the vector $\Vc_e(t)$ defined by~\eqref{VW} is a solution of the ordinary differential equation
\begin{equation}\label{deqv2}
\nabla_{T_e}\Vc_e=-\frac{1}{\|T_e\|_{G_0}}R\left(V_e,T_e\right)T_e.
\end{equation}
with initial vector $\Vc_e(h_{0,\,e}(0))$ if $\|T_e\|_{G_0}\neq 0$.
\end{rema}




We postpone the proof of Proposition~\ref{key2} for the moment. Then, for any Weil--Petersson geodesic $\{[G_s]\}_{s\,\in \,(-\e,\,\e)}$ in $\Tc(S)$, the second variation formula~\eqref{esvf2} becomes
\begin{multline}\label{svff}
\frac{d^2}{d s^2}\Ec_{\Cc}\left([G_{s}]\right)\bigg|_{s=0}\\
=\sum_{e\,\in\,E}m_E(e)\int_0^1\left\{-\frac{1}{4}\left(\Delta_{G_0}-2\right)^{-1}\left\|Q\right\|_{G_0}^2+\left\|V_e^{\perp}\right\|_{G_0}^2+\left\|\Vc_e\right\|_{G_0}^2\right\}\cdot \left\|T_e\right\|^2\,dt,
\end{multline}
by Propositions~\ref{key1} and~\ref{key2}, and using the fact
\[
\frac{1}{2}\sum_{e\,\in\,E}m_E(e)\int_0^1 h_0^*L_ZG\left(T_e,T_e\right)dt=0
\]
due to the harmonicity of $h_0$ (see~\eqref{harm} for a similar identity). More generally, we obtain the following Hessian formula:

\begin{theo}\label{Thm:Hessian}
Let $\{[G_{r,\,s}]\}_{r,\,s\,\in\,(-\e,\,\e)}$ be any two-parameter family of Weil--Petersson geodesics in $\Tc(S)$ with $[G_0]=[G_{0,\,0}]$ with a horizontal lift $\{G_{r,\,s}\}_{r,\,s\,\in\,(-\e,\,\e)}$ in $\Mc_{-1}(S)$ for $\e>0$. For the tangent vectors $P:=(\p_r G_{r,\,0})|_{r=0}$ and $Q:=(\p_s G_{0,\,s})|_{s=0}$, the Hessian of the energy functional $\Ec_{\Cc}: \Tc(S)\to \R$ at $[G_0]$ is given by
\begin{multline*}
\Hess_{\Ec_\Cc}\left(P,Q\right)_{[G_0]}=\frac{\p^2}{\p r\p s}\Ec_{\Cc}\left([G_{r,\,s}]\right)\Big|_{r=s=0}=\\
\sum_{e\,\in\,E}m_E(e)\int_0^1\left\{-\frac{1}{4}\left(\Delta_{G_0}-2\right)^{-1}\langle P, Q\rangle_{G_0} +\left\langle V_e^{\perp}, W_e^{\perp}\right\rangle_{G_0}+\left\langle \Vc_e, \Wc_e\right\rangle_{G_0}\right\}\cdot \|T_e\|^2_{G_0}dt,
\end{multline*}
where $h_{0}:X \to (S, G_{0})$ is any smooth harmonic map, $T_e(t):=\p_th_{0,\,e}(t)$, $V_e(t):=\p_r h_{r,\,e}(t)$ for the family of harmonic maps $h_r: X\to (S, G_{r,\,0})$, $\Vc_e(t)$ are defined by~\eqref{VW} relative to the variable $r$, and $W_e(t)$, $\Wc_e(t)$ are defined analogously relative to the variable $s$. In the integral, $\langle P, Q\rangle_{G_0}$ stands for $G_0(P, Q)$.
\end{theo}

One can check this formula by using~\eqref{svff}, and thus, we omit the proof. Now, we give a proof of strictly convexity of the energy functional $\Ec_{\Cc}$.



\begin{proof}[{Proof of Theorem~\ref{Thm:convex}}]
We show that
\begin{equation}\label{Lem: Wolf}
-\frac{1}{4}\left(\Delta_{G_0}-2\right)^{-1}\|Q\|^2_{G_0}\geq \frac{1}{24}\|Q\|^2_{G_0}\geq 0.
\end{equation}
Indeed, setting $u:=(\Delta_{G_0}-2)^{-1}\|Q\|_{G_0}^2+\tfrac{1}{6}\|Q\|_{G_0}^2$, we see
\[
\left(\Delta_{G_0}-2\right)u=\|Q\|_{G_0}^2+\frac{1}{6}\left(\Delta_{G_0}-2\right)\|Q\|_{G_0}^2=\frac{1}{6}\left(\Delta_{G_0}+4\right)\|Q\|_{G_0}^2\geq 0,
\]
where the last inequality follows from Lemma~\ref{Lem:Qholomorphic}$\MK$\eqref{lemm3.1.2}. Suppose ${\max}_{p\,\in\,S}\,u(p)>0$. Then, at any maximum point $p_0$ of $u$, we have $\Delta_{G_0} u(p_0)<0$ and $-2u(p_0)<0$, which leads $(\Delta_{G_0}-2)u(p_0)<0$, a contradiction. Hence $u\leq 0$ on $S$, as required.



Combining~\eqref{Lem: Wolf} with~\eqref{svff}, we have
\begin{multline}\label{hesq}
\Hess_{\Ec_{\Cc}}(Q,Q)_{G_0}\\
\geq \sum_{e\,\in\,E}m_E(e)\int_0^1\left\{\frac{1}{24}\|Q\|^2_{G_0}+\left\|V_e^{\perp}\right\|_{G_0}^2+\left\|\Vc_e\right\|_{G_0}^2\right\}\cdot \left\|T_e\right\|_{G_0}^2\, dt\geq 0
\end{multline}
for any Weil--Petersson geodesic $\{[G_s]\}_{s\,\in\,(-\e,\,\e)}$ in $\Tc(S)$, and thus $\Ec_{\Cc}$ is convex.


Suppose there exists $Q$ such that $\Hess_{\Ec_{\Cc}}(Q,Q)_{G_0}=0$. Then~\eqref{hesq} implies
\begin{equation}\label{hol}
Q(h_{0,\,e}(t))\equiv 0
\end{equation}
for any edge $e\in E$ with $\|T_e\|_{G_0}\neq 0$. For such an edge $e\in E$, we take a point $p$ on the image of $h_e$ and an isothermal coordinate $z=x^1+ix^2$ around $p$. Recall that, by Lemma~\ref{Lem:Qholomorphic}, the function $Q_{11}-\sqrt{-1}Q_{12}$ is a holomorphic function for $Q\in T_{[G]}\Tc(S)$. Since the zeros of a holomorphic function are isolated, \eqref{hol} implies $Q\equiv 0$ on an open subset around the point $p$. Then, it turns out that $Q\equiv 0$ on the whole $S$ because $S$ is a (compact) connected surface. Therefore for any nonzero $Q\in T_{[G]}\Tc(S)$, we have $\Hess_{\Ec_{\Cc}}(Q,Q)_{G_0}>0$, and thus $\Ec_{\Cc}$ is strictly convex.
\end{proof}

\begin{rema}
In fact, the above proof shows that in the setting of Theorem~\ref{Thm:convex}, the Hessian of the energy functional $\Ec_{\Cc}: \Tc(S)\to \R$ is positive definite and non-degenerate at any point $[G]\in \Tc(S)$.
\end{rema}



\subsection{Proof of Proposition~\ref{key2}} Let $\{[G_s]\}_{s\,\in\,(-\e,\,\e)}$ be a smooth curve in $\Tc(S)$, and $h_{s}$ the family of harmonic maps associated with its horizontal lift $\{G_s\}_{s\,\in\,(-\e,\,\e)}$ in $\Mc_{-1}(S)$. For the simplicity of notation, let us write
\[
T_e^{s}(t):=\p_t h_{s,\,e}(t), \quad \text{and} \quad V_e^{s}(t):=\p_s h_{s,\,e}(t).
\]
for each edge $e\in E$. Moreover, we simply write $T_e(t):=T_e^{0}(t)$ at $s=0$. Since $h_{s}: X\to (S, G_{s})$ is a family of harmonic maps, we may assume the following:
\begin{itemize}
\item $h_{s,\,e}(t)$ is a constant speed geodesic in $(S, G_{s})$ for each $e\in E$, namely, we have
\begin{align}\label{as21}
\nabla^{s}_{T_e^{s}}T_e^{s}=0
\end{align}
along $h_{s,\,e}(t)$ for any $s \in (-\epsilon, \epsilon)$, where $\nabla^{s}$ is the Levi--Civita connection of $(S, G_{s})$. We simply write $\nabla=\nabla^0$. The parameter $t$ is proportional to the arc-length parameter of $h_{s,\,e}(t)$, i.e., the norm $\|T_e^{s}(t)\|_{G_{s}}$ is constant along $h_{s,\,e}(t)$ in $t \in [0, 1]$ for each $e\in E$ and each $s \in (-\epsilon, \epsilon)$.
\item The balanced condition holds for the map $h_{s}: X\to (S, G_{s})$,
\begin{equation}\label{as22}
\sum_{e\,\in\,E_x}m_E(e) T_e^{s}(0)=0
\end{equation}
for any $x\in V$ and all $s\in (-\epsilon, \epsilon)$.
\end{itemize}



The second variation formula~\eqref{Eq:secondVF} in Appendix implies that
\begin{equation}\label{hf1}
\frac{d^2}{d s^2} E_{G_0}(h_{s})\Big|_{s=0}=\sum_{e\,\in\,E}m_E(e) \int_0^1 \Big\{\left\|\nabla_{T_e}V_e\right\|_{G_0}^2 -G_0\left(R\left(V_e, T_e\right)T_e, V_e\right)\Big\}\,dt
\end{equation}
for the harmonic map $h_{0}: X\to (S,G_0)$, where $R$ is the curvature tensor of $(S, G_0)$. Then, the right hand side becomes
\begin{multline}\label{hf2}
\sum_{e\,\in\,E}m_E(e)\int_0^1\Big\{-G_0\left(\nabla_{T_e}\nabla_{T_e}V_e, V_e\right)-G_0\left(R\left(V_e, T_e\right)T_e, V_e\right)\Big\}\,dt\\
+\sum_{e\,\in\,E}m_E(e)\big[G_0\left(\nabla_{T_e}V_e, V_e\right)\big]_{t=0}^{t=1}.
\end{multline}
The following lemma is the place where we use~\eqref{as22} that the balanced condition holds for all $s \in (-\e, \e)$.


\begin{lemm}\label{le2}
Letting $T_e^{s}(t):=\p_t h_{s,\,e}(t)$ and $V_e^{s}(t):=\p_s h_{s,\,e}(t)$, if we have
\begin{equation}\label{Eq:BCall-s}
\sum_{e\,\in\,E_x}m_E(e)T_e^s(0)=0 \quad \text{for each }x \in V\text{ for all }s \in (-\e, \e),
\end{equation}
then $\nabla_{V_x}\cdot (\sum_{e\,\in\,E_x}m_E(e)T_e)=0$ at each $x \in V$. In particular, it holds that
\[
\sum_{e\,\in\,E}m_E(e)\big[G_0\left(\nabla_{T_e}V_e, V_e\right)\big]_{t=0}^{t=1}=0.
\]
\end{lemm}


\begin{proof}
For each vertex $x$, differentiating~\eqref{Eq:BCall-s} in $s \in (-\e, \e)$, we have
\[
\sum_{e\,\in\,E_x}m_E(e)\nabla_{V_x}T_e=\nabla_{V_x}\cdot\left(\sum_{e\,\in\,E_x}m_E(e)T_e\right)=0.
\]
Denote the origin and the terminal of edge $e\in E$ by $o(e)$ and $t(e)$, respectively. Since
\begin{gather*}
G_0\left(V_e, \nabla_{T_e}V_e\right)_{h_0(t(e))} =-G_0\left(V_{\overline{e}}, \nabla_{T_{\overline{e}}}V_{\overline{e}}\right)_{h_0\left(o\left(\overline{e}\right)\right)}
\\
\intertext{for}
 T_{\overline{e}} =-T_e\quad\text{and}\quad m_E\left(\overline{e}\right)=m_E(e),
\end{gather*}
we have that
\begin{equation*}\label{e1}
\sum_{e\,\in\,E}m_E(e)\Big[G_0\left(V_e, \nabla_{T_e}V_e\right)\Big]_{t=0}^{t=1} =-2\sum_{e\,\in\,E}m_E(e)G_0\left(V_e, \nabla_{T_e}V_e\right)_{h_0(o(e))}.
\end{equation*}
Note that $[V_e, T_e]=0$ and thus $\nabla_{T_e}V_e=\nabla_{V_e}T_e$. For each vertex $x$, letting $V_x=V_e$ for $e \in E_x$ we have
\begin{align*}
-2\sum_{x\,\in\,V}\,\sum_{e\,\in\,E_x}m_E(e)G_0\left(V_e, \nabla_{T_e}V_e\right)_{h_0(x)} &=-2\sum_{x\,\in\,V}G_0\left(V_x, \sum_{e\,\in\,E_x}m_E(e)\nabla_{V_x}T_e\right)_{h_0(x)}\\
&=0,
\end{align*}
as required.
\end{proof}



Recall that, for the velocity vector $Q:=(\p_sG_s)|_{s=0}$, we have defined the vector field $Q_e^{\sharp}$ along $h_{0,\,e}(t)$ such that $Q(T_e, \cdot)=G_0(Q_e^{\sharp},\cdot)$ (Definition~\ref{Def:V}). The following formula uses the property~\eqref{as21} that $h_{s,\,e}(t)$ is a constant speed geodesic in $(S, G_s)$ for each $e \in E$ for all $s \in (-\e, \e)$.

\begin{lemm}\label{le3}
For any horizontal lift $\{G_s\}_{s\,\in\,(-\e,\,\e)}$ of any smooth curve \linebreak $\{[G_s]\}_{s\,\in\,(-\e,\,\e)}$ for $\e>0$ in $\Tc(S)$, we have
\[
-\nabla_{T_e}\nabla_{T_e}V_e-R\left(V_e, T_e\right)T_e=\frac{1}{2}\nabla_{T_e}Q_{e}^{\sharp} \quad \text{for }\; e \in E.
\]
\end{lemm}

\begin{proof}
We shall consider the differential of~\eqref{as21} at $s=0$. For any fixed $t \in [0, 1]$, we may take a local isothermal coordinate $(x^1,x^2)$ around $h_{0,\,e}(t)$ satisfying $G_0(h_{0,\,e}(t))=4((dx^1)^2+(dx^2)^2)$ and $\Gamma_{\beta\gamma}^{\alpha}(G_0, h_{0,\,e}(t))=0$ at the fixed point $h_{0,\,e}(t)$. Indeed, since $(S,G_0)$ is a hyperbolic surface, the fixed point $h_{0,\,e}(t)$ can be regarded as the origin of the Poincaré disk, and the standard coordinate $(x^1,x^2)$ of the Poincaré disk gives rise to the desired local coordinate.

Let
\[
T_e^s(t)=\p_t h_{s,\,e}(t)=\tau^\a_s(t)\frac{\p}{\p x^\a}\Big|_{h_{0,\,e}(t)} \qquad \text{for}\;\a=1, 2.
\]
Then, we have for $\alpha=1,2$,
\begin{multline}\label{c4}
\frac{\p}{\p s}\Big|_{s=0}\left(\nabla^s_{T_e^s}T_e^s\right)^{\alpha} =\frac{\p}{\p s}\Big|_{s=0}\left(\p_t \t_s^{\alpha}+\t_s^{\beta}\t_s^{\gamma}\Gamma_{\beta\gamma}^{\alpha}\left(G_s, h_s\right)\right)\\
\begin{aligned}
&=\frac{\p}{\p s}\Big|_{s=0}\left(\p_t \t_s^{\alpha}\right)+\left(\frac{\p}{\p s}\middle|_{s=0}\t_s^{\beta}\t_s^{\gamma}\right)\Gamma_{\beta\gamma}^{\alpha}\left(G_0, h_0\right)\\
&\mkern124mu +\t_0^{\beta}\t_0^{\gamma}\left(\frac{\p}{\p s}\middle|_{s=0}\Gamma_{\beta\gamma}^{\alpha}\left(G_s, h_0\right)+\frac{\p}{\p s}\middle|_{s=0}\Gamma_{\beta\gamma}^{\alpha}\left(G_0, h_s\right)\right)\\
&=\frac{\p}{\p s}\Big|_{s=0}\left(\p_t \t_s^{\alpha}+\t_s^{\beta}\t_s^{\gamma}\Gamma_{\beta\gamma}^{\alpha}\left(G_0,h_s\right)\right)+\frac{\p}{\p s}\Big|_{s=0}\left(\p_t \t_0^{\alpha}+\t_0^{\beta}\t_0^{\gamma}\Gamma_{\beta\gamma}^{\alpha}\left(G_s, h_0\right)\right)\\
&=\frac{\p}{\p s}\Big|_{s=0}\left(\nabla^0_{T_e^s}T_e^s\right)^{\alpha}+\frac{\p}{\p s}\Big|_{s=0}\left(\nabla^s_{T_e}T_e\right)^{\alpha},
\end{aligned}
\end{multline}
where we have added $\frac{\p}{\p s}|_{s=0}(\p_t \t_0^{\alpha})=0$ in the third equality, $\Gamma_{\beta\gamma}^{\alpha}(G_s, h_s)$ denotes the Christoffel symbol for the metric $G_s$ at the point $h_{s,\,e}(t)$ (and we use the Einstein convention for summations over $\b, \g=1, 2$). Since $\Gamma_{\beta\gamma}^{\alpha}(G_0, h_{0,\,e}(t))=0$, we have
\[
\frac{\p}{\p s}\Big|_{s=0}\left(\nabla^0_{T_e^s}T_e^s\right)^{\alpha}=\left(\nabla_{V_e^s}\nabla_{T_e^s}T_e^s\right)^\a|_{s=0},
\]
and since $[V_e^s, T_e^s]=0$,
\begin{align*}
\left(\nabla_{V_e^s}\nabla_{T_e^s}T_e^s\right)\big|_{s=0} &=\left(R\left(V_e^s, T_e^s\right)T_e^s+\nabla_{T_e^s}\nabla_{V_e^s}T_e^s\right)\big|_{s=0}\\
&=\left(R\left(V_e^s, T_e^s\right)T_e^s+\nabla_{T_e^s}\nabla_{T_e^s}V_e^s\right)\big|_{s=0}
\end{align*}
Therefore we obtain for $\a=1, 2$,
\begin{equation}\label{c41}
\frac{\p}{\p s}\Big|_{s=0}\left(\nabla^0_{T_e^s}T_e^s\right)^{\alpha}=\big(R\left(V_e, T_e\right)T_e+\nabla_{T_e}\nabla_{T_e}V_e\big)^\a.
\end{equation}



Consider the second term in~\eqref{c4}. We simply write $\t^i=\t^i_0$ and $\Gamma_{\beta\gamma}^{\alpha}(s)=\Gamma_{\beta\gamma}^{\alpha}(G_s, h_0)$ so that
\[
\nabla^s_{T_e}T_e=\left\{\p_t\tau^\a+ \tau^\b \tau^\g\Gamma_{\b \g}^\a(s)\right\}\frac{\p}{\p x^\a}\Big|_{h_{0,\,e}(t)}
\]
Since $\nabla_{T_e}T_e$ and $\Gamma_{\b \g}^\a(0)=0$ at $t$, we have $\p_t\tau^\a(t)=0$ for $\a=1,2$. Concerning the differential of Christoffel symbol $\Gamma_{ij}^k(s)$ with respect to $s$, we have
\begin{align*}
\frac{\p}{\p s}\Gamma_{ij}^k(s)\Big|_{s=0}&=\frac{\p}{\p s}\frac{1}{2}(G_s)^{km}\Big\{(G_s)_{im,\,j}+(G_s)_{jm,\,i}-(G_s)_{ij,\,m}\Big\}\Big|_{s=0}\\
&=\frac{1}{8}\Big\{\left(\p_s G_s\right)_{ik,\,j}+\left(\p_s G_s\right)_{jk,\,i}-\left(\p_s G_s\right)_{ij,\,k}\Big\}\Big|_{s=0}\
\end{align*}
at the point $h_{0,\,e}(t)$ since $(G_0)^{km}(t)=(1/4)\delta_{km}$ and $(G_0)_{ij,\,k}(t)=0$ for any $i,j,k=1, 2$. Recall that $\{G_s\}_{s\,\in\,(-\e,\,\e)}$ is a horizontal lift of a curve $\{[G_s]\}_{s\,\in\,(-\e, \,\e)}$ in $\Tc(S)$; Lemma~\ref{Lem:Qholomorphic} implies that $(\p_s G_s)_{11}-\sqrt{-1}(\p_s G_s)_{12}$ satisfies the Cauchy--Riemann equation around $h_{0,\,e}(t)$ and that $(\p_s G_s)_{ij,\,k}|_{s=0}$ is symmetric under permuting any indices $i, j, k \in \{1, 2\}$. Hence we have
\[
\frac{\p}{\p s}\Gamma_{ij}^k(s)\Big|_{s=0} =\frac{1}{8}\frac{\p}{\p s}({G}_s)_{ik,\,j}\Big|_{s=0} \quad \text{at }h_{0,\,e}(t).
\]
Therefore we see that
\begin{align*}
\frac{\p}{\p s}\Big|_{s=0}\left(\nabla^s_{T_e}T_e\right)^{\alpha} &=\frac{\p}{\p s}\Big\{\p_t\tau^\a+\tau^\b \tau^\g\Gamma_{\b \g}^\a(s)\Big\}\Big|_{s=0}\\
&=\tau^\b \tau^\g \left(\frac{\p}{\p s}\Gamma_{\b \g}^\a(s)\right)\Big|_{s=0} =\frac{1}{8}\tau^\b \tau^\g\frac{\p}{\p s}\left({G}_s\right)_{\b\a,\,\g}\Big|_{s=0},
\end{align*}
and thus
\begin{multline}\label{c42}
\frac{\p}{\p s}\Big|_{s=0}\left(\nabla^s_{T_e}T_e\right)^{\alpha}\\
\begin{aligned}
&=\frac{1}{8}\frac{\p}{\p t}\left\{\frac{\p}{\p s}{G}_s\middle|_{s=0}\left(T_e, \frac{\p}{\p x^\a}\right)\right\}=\frac{1}{8}\frac{\p}{\p t}Q_e\left(T_e, \frac{\p}{\p x^\a}\right)\\
&=\frac{1}{8}\frac{\p}{\p t}G_0\left(Q_e^{\sharp}, \frac{\p}{\p x^\a}\right)=\frac{1}{8}G_0\left(\nabla_{T_e} Q_e^{\sharp}, \frac{\p}{\p x^\a}\right)=\frac{1}{2}\left(\nabla_{T_e}Q_e^{\sharp}\right)^{\alpha}
\end{aligned}
\end{multline}
where $(\nabla_{T_e}\frac{\p}{\p x^\a})(h_{0,\, e}(t))=0$ and $G_0(h_{0,\,e}(t))=4((dx^1)^2+(dx^2)^2)$ in this coordinate.


Finally, substituting~\eqref{c41} and~\eqref{c42} to~\eqref{c4}, we obtain for each $t \in [0, 1]$,
\[
\frac{\p}{\p s}\Big|_{s=0}\left(\nabla^s_{T_e^s}T_e^s\right)=\nabla_{T_e}\nabla_{T_e}V_e+R\left(V_e, T_e\right)T_e+\frac{1}{2}\nabla_{T_e}Q_e^{\sharp},
\]
and since we have~\eqref{as21}, we obtain the Lemma~\ref{le3}.
\end{proof}



\begin{proof}[{Proof of Proposition~\ref{key2}}]
By Lemma~\ref{le2} and Lemma~\ref{le3}, the consequence of the second variation formula~\eqref{hf2} becomes
\begin{multline*}
\frac{d^2}{d s^2} E_{G_0}(h_{s})\Big|_{s=0}\\
\begin{aligned}
&=\sum_{e\,\in\,E}m_E(e)\int_0^1 \frac{1}{2}G_0\left(\nabla_{T_e}Q_{e}^{\sharp},V_e\right) \,dt\\
&= \frac{1}{2}\sum_{e\,\in\,E}m_E(e)\int_0^1-G_0\left(Q_e^{\sharp},\nabla_{T_e}V_e\right)\,dt+\frac{1}{2}\sum_{e\,\in\,E}m_E(e)\Big[Q\left(T_e,V_e\right)\Big]_{t=0}^{t=1}.
\end{aligned}
\end{multline*}
The second term is $0$ since the balanced condition~\eqref{as22} implies that
\begin{align*}
\frac{1}{2}\sum_{e\,\in\,E}m_E(e)\Big[Q\left(T_e,V_e\right)\Big]_{t=0}^{t=1} &=-\sum_{e\,\in\,E}m_E(e)Q\left({T_e}(0), V_e(0)\right)\\
&=-Q\left(\sum_{e\,\in\,E}m_E(e)T_e(0), V_e(0)\right)=0.
\end{align*}
The integrand of the first term for each $e \in E$ is
\begin{equation}\label{svf3}
-G_0\left(Q_e^{\sharp}, \nabla_{T_e}V_e\right) =\left\|\nabla_{T_e}V_e\right\|_{G_0}^2+\frac{\left\|Q_e^{\sharp}\right\|_{G_0}^2}{4}-\left\|\nabla_{T_e}V_e+\frac{Q_e^{\sharp}}{2}\right\|^2_{G_0}.
\end{equation}
Letting $J_0$ be the complex structure on $(S, G_0)$, we have for $e \in E$ with $T_e \neq 0$,
\[
\left\|Q_e^{\sharp}\right\|_{G_0}^2=\frac{1}{{\left\|T_e\right\|_{G_0}^2}}\left\{Q \left(T_e,T_e\right)^2+Q\left(T_e, J_0 \cdot T_e\right)^2\right\} =\frac{\|Q\|_{G_0}^2\|T_e\|_{G_0}^2}{2},
\]
and note that by Definition~\ref{Def:V} for $e \in E$ with $T_e \neq 0$,
\[
\Vc_e=\frac{1}{\|T_e\|_{G_0}}\left(\nabla_{T_e}V_e+\frac{Q_e^{\sharp}}{2}\right),
\]
and by~\eqref{hf2},
\begin{multline*}
\sum_{e\,\in\,E}m_E(e)\int_0^1\|\nabla_{T_e}V_e\|_{G_0}^2\, dt \\
=\frac{d^2}{ds^2}E_{G_0}(h_s)\Big|_{s=0}+\sum_{e\,\in \,E}m_E(e)\int_0^1 G_0\left(R\left(V_e, T_e\right)T_e,V_e\right)\, dt.
\end{multline*}
Substituting these to~\eqref{svf3}, we obtain
\begin{multline*}
\frac{1}{2}\cdot \frac{d^2}{d s^2} E_{G_0}(h_{s})\Big{|}_{s=0}\\
=\frac{1}{2}\sum_{e\,\in\,E}m_E(e)\int_0^1\left\{G_0\big(R\left(V_e, T_e\right)T_e,V_e\big)+\frac{\|Q\|_{G_0}^2\|T_e\|_{G_0}^2}{8}-\|\Vc_e\|_{G_0}^2 \|T_e\|_{G_0}^2\right\}\, dt.
\end{multline*}
Since $(S,G_0)$ is a surface of constant sectional curvature $-1$, we have
\[
G_0\big(R(V_e, T_e)T_e,V_e\big)=-\left\|V_e^{\perp}\right\|^2_{G_0}\|T_e\|_{G_0}^2,
\]
where $V_e^{\perp}$ is the orthonormal projection with respect to $T_e$ for $T_e=0$ and $0$ otherwise, and we complete the proof of Proposition~\ref{key2}.
\end{proof}



\section{Properness of the energy functional}\label{Sec:proper}

We will establish the properness of the energy functional $\Ec_\Cc$ for $\Cc=[f]$ if $f:X \to S$ induces a surjective homomorphism between the fundamental groups. First we show the following lemma (which will also be used in the next Section~5).


\begin{lemm}\label{Lem:homotopic}
Assume that $f:X \to S$ is a continuous map which induces a surjective homomorphism from $\pi_1(X)$ to $\pi_1(S)$. If $\f_1, \f_2$ are orientation-preserving homeomorphisms on $S$ and $\f_1 \circ f$ and $\f_2 \circ f$ are homotopic, then $\f_1$ and $\f_2$ are isotopic on $S$.
\end{lemm}

\begin{proof}
Let $\wh \f:=\f_2^{-1}\circ \f_1$, then the maps $\wh \f \circ f$ and $f$ are homotopic. For any $x_0 \in X$, let $\g$ be the path from $f(x_0)$ to $\wh \f \circ f(x_0)$ given by the homotopy between $\wh \f \circ f$ and $f$. Then, there exists a homeomorphism $\b_\g$ homotopic to the identity map $\id$ on $S$ such that
\[
\b_{\g \ast}:\pi_1\left(S, \wh \f \circ f(x_0)\right) \to \pi_1(S, f(x_0)), \qquad \b_{\g \ast}([\a])=[\g \cdot \a \cdot \wbar \g],
\]
for $[\a] \in \pi_1(S, \wh \f \circ f(x_0))$, where $\wbar \g$ denotes the reversed path of $\g$. (Such an $\b_\g$ is obtained first when $\g$ is a short enough path on $S$, and in general divide $\g$ into finitely many short enough paths and composite homeomorphisms homotopic to the identity.) Since $f$ induces a surjective homomorphism $f_\ast: \pi_1(X, x_0) \to \pi_1(S, f(x_0))$, one obtains $\b_{\g \ast} \circ \wh \f_\ast: \pi_1(S, f(x_0)) \to \pi_1(S, f(x_0))$,
\[
[f(c)] \mapsto \left[\g \cdot \wh \f (f(c)) \cdot \wbar \g\right] \quad \text{for }[c] \in \pi_1(X, x_0),
\]
and since $\wh \f \circ f$ and $f$ are homotopic, $\b_{\g \ast} \circ \wh \f_\ast$ coincides with $\id_\ast$ given by $\id$ on $S$. Since $S$ is a $K(\pi_1(S), 1)$-space, $\b_\g \circ \widehat \f$ and $\id$ are homotopic fixing $f(x_0)$~\cite[Proposition~1B.9]{Hatcher}, and since $\b_\g$ is homotopic to $\id$, $\widehat \f$ is homotopic to $\id$; hence $\f_1$ and $\f_2$ are homotopic. It is known that two orientation-preserving homeomorphisms which are homotopic on a closed Riemann surface $S$ are isotopic~\cite[Theorem~6.3]{Epstein-isotopies}, therefore $\f_1$ and $\f_2$ are in fact isotopic on $S$.
\end{proof}





\begin{proof}[{Proof of Theorem~\ref{Thm:proper}}]
For any $R \ge 0$, let us denote the sublevel set by
\[
\Sub(R):=\Big\{\big[(\SS, G_\SS), \f\big] \in \Tc(S) \ : \ \Ec_\Cc([G]) \le R, \, G=\f^\ast G_\SS\Big\}.
\]
Assume that $\Sub(R) \neq \emptyset$. Taking a harmonic map $h_G: X \to (S, G)$ in the homotopy class $\Cc=[f]$, we write
\[
\Ec_\Cc([G])=E_G(h_G)=\frac{1}{2}\sum_{e\,\in \,E}m_E(e)\int_0^1\left\|\partial_t h_{G,\,e}\right\|_G^2\, dt.
\]
Let $l(h_{G,\,e})$ be the length of $h_{G, \,e}:[0, 1] \to (S, G)$ for each $e \in E$. Since $h_G$ is harmonic and piecewise geodesic, we have for each $e \in E$,
\[
l\left(h_{G,\,e}\right)=\int_0^1\left\|\partial_t h_{G,\,e}\right\|_G\, dt=\left\|\partial_t h_{G,\,e}\right\|_G.
\]
The Cauchy--Schwarz inequality gives
\begin{align*}
\sum_{e\,\in\,E}l\left(h_{G,\,e}\right) &= \sum_{e\,\in\,E}\left(\int_0^1 \left\|\partial_t h_{G,\,e}\right\|_G^2 \, dt\right)^{1/2} \\
&\le \left(\sum_{e\,\in\,E}\frac{1}{m_E(e)}\right)^{1/2}\cdot \left(\sum_{e\,\in\,E}m_E(e)\int_0^1 \left\|\partial_t h_{G,\,e}\right\|_G^2 \, dt\right)^{1/2}.
\end{align*}
Hence letting $M:=\sum_{e\,\in\,E}m_E(e)^{-1}$, we obtain for any $[G] \in \Sub(R)$,
\begin{equation}\label{Eq:length}
\sum_{e\,\in\,E}l\left(h_{G,\,e}\right) \le \left(M\cdot 2\Ec_{\Cc}([G])\right)^{1/2} \le (2 MR)^{1/2}.
\end{equation}


For given closed Riemann surface $S$, we take a finite collection of simple closed curves $\{\g_1,\,\dots,\,\g_N\}$ on $S$ such that the union $\g_1 \cup\,\cdots\,\cup \g_N$ fills $S$, i.e., the complement of this union is a disjoint union of disks. (For example, one may take simple closed curves corresponding to the free homotopy classes of the standard set of generators of $\pi_1(S)$, where $N$ is twice the genus of $S$.) Then, if $\g$ is any simple closed curve which is not homotopic to a point, then there exists a $\g_i$ such that the geometric intersection number between free homotopy classes of $\g$ and $\g_i$ is not zero.



Since $f_\ast: \pi_1(X, x_0) \to \pi_1(S, f(x_0))$ is surjective, we choose a collection of closed paths $\{c_1,\,\dots,\, c_N\}$ in $X$ such that $f(c_i)$ is freely homotopic to $\g_i$ for each $i=1,\,\dots,\,N$. Let $C_{\max}:=\max_{i=1,\,\dots,\,N} |c_i|$, where $|c_i|$ is the number of edges in $c_i$. (One may just take $C_{\max}=1$ if $f$ fills the surface in the following inequality.) Then, for any $c_i$, we have
\[
\sum_{e\,\in\,c_i}l\left(h_{G,\,e}\right) \le C_{\max}\, \sum_{e\,\in\,E}l\left(h_{G,\,e}\right) \le C_{\max}\,(2 MR)^{1/2},
\]
for any $[G] \in \Sub(R)$. For any closed hyperbolic surface $(S, G)$, let $\inj((S, G))$ be the injectivity radius. Taking a closed geodesic $\g_{\inj}$ in $(S, G)$ such that the length realizes twice the injectivity radius $l(\g_{\inj})=2\, \inj ((S, G))$, we have a $\g_i$ such that the geometric intersection number is non zero in their free homotopy classes. Let $\g_{i,\,G}$ be a (unique) closed geodesic freely homotopic to $\g_i$, then $\g_{i,\,G}$ has the shortest length among the free homotopy class of $\g_i$ since $(S, G)$ is a hyperbolic surface, and we have
\[
l(\g_{i,\,G}) \le \sum_{e\,\in\,c_i}l(h_{G, \,e}) \le C_{\max}\,(2 MR)^{1/2}.
\]
Then, the Collar Lemma~\cite{Keen} (e.g., \cite[Lemma~13.6]{FarbMargalit}) implies that
\[
N_{\g_i}=\left\{x \in S \, : \, d(x, \g_i) \le \sinh^{-1}\left(\frac{1}{\sinh (l(\g_i)/2)}\right)\right\}
\]
is an embedded annulus, and thus
\[
l(\g_{\inj}) \ge \sinh^{-1}\left(\frac{1}{\sinh (l(\g_i)/2)}\right).
\]
Therefore, letting
\[
\e(R):=\frac{1}{2}\,\sinh^{-1}\left(\frac{1}{\sinh ((1/2)C_{\max}\,(2 MR)^{1/2})}\right)>0,
\]
we obtain $\inj ((S, G)) \ge \e(R)$ for any $[G] \in \Sub(R)$. We consider the moduli space $\Mc(S)$ of $S$, and define
\[
\Mc_{\e(R)}(S):=\big\{[(S, G)] \in \Mc(S) \, : \, \inj((S, G))\ge \e(R) \big\}.
\]
Now we are discussing the standard topology in Teichm\"uller space $\Tc(S)$, and the induced topology in the moduli space $\Mc(S)$ as the quotient space of $\Tc(S)$ by the action of the mapping class group $\Mod(S)$ which acts on $\Tc(S)$ properly discontinuously. Then, the Mumford--Mahler Compactness criterion implies that $\Mc_{\e(R)}(S)$ is compact (\cite[Theorem~C.1]{TrombaBook}); for any sequence of points $[(\SS_i, G_i), \f_i]$ in $\Tc(S)$ with $\Ec_\Cc([\f_i^\ast G_i]) \le R$, we have $\inj((S, \f_i^\ast G_i)) \ge \e(R)$, and thus up to passing to a subsequence there exists a hyperbolic surface $[(S, G_\star)]$ in $\Mc_{\e(R)}(S)$ such that $[(S, \f_i^\ast G_i)]$ converges to $[(S, G_\star)]$ in $\Mc(S)$. Moreover, there exists a sequence of (orientation-preserving) diffeomorphisms $\wh \f_i$ on $S$ to itself and $\f_\star:=\id:S \to S$ such that $[(\SS_i, G_i), \f_i \circ \wh \f_i^{-1}]$ converges to $[(S, G_\star), \f_\star]$ in the Teichm\"uller space $\Tc(S)$, and $(\f_i \circ \wh \f_i^{-1})^\ast G_i$ converges to $G_\star$ in the $C^\infty$-topology: for any $\e>0$, for all large enough~$i$, we have
\begin{equation}\label{Eq:dilatation}
G_\star \le (1+\e)\left(\f_i \circ \wh \f_i^{-1}\right)^\ast G_i.
\end{equation}

Let $h_i:X \to (S, \f_i^\ast G_i)$ be a harmonic map in the homotopy class $\Cc$. The maps $\wh \f_i \circ h_i:X \to (S, (\f_i \circ \wh \f_i^{-1})^\ast G_i)$ are harmonic since $\wh \f_i: (S, \f_i^\ast G_i) \to (S, (\f_i \circ \wh \f_i^{-1})^\ast G_i)$ are isometry (although $\wh \f_i \circ h_i$ may not be homotopic to $h_i$). We note that $\wh \f_i \circ h_i: X \to (S, G_\star)$ are equicontinuous; this follows from~\eqref{Eq:length} and~\eqref{Eq:dilatation}. Thus by the Ascoli--Arzel\`a theorem after passing to a subsequence if necessary, $\wh \f_i \circ h_i$ converges to a continuous map $h_\star: X \to (S, G_\star)$ uniformly. Therefore for all large enough $i, j$, the maps $\wh \f_i \circ h_i$ and $\wh \f_j \circ h_j$ are homotopic.

Then, the condition that $\wh \f_i \circ h_i$ and $\wh \f_j \circ h_j$ are in the same homotopy class inducing surjective homomorphisms $\pi_1(X) \to \pi_1(S)$ implies that $\wh \f_i$ and $\wh \f_j$ are isotopic on $S$. Indeed, since both $h_i$ and $h_j$ are homotopic to $f$, $\wh \f_i \circ f$ and $\wh \f_j \circ f$ are homotopic for all large enough $i, j$, and Lemma~\ref{Lem:homotopic} implies that $\wh \f_i$ and $\wh \f_j$ are isotopic for all large enough $i, j$.

Let $\wh \f_\star:=\wh \f_i$ for a large enough $i$. Since $[(\SS_i, G_i), \f_i \circ \wh \f_i^{-1}]$ converges to $[(S, G_\star), \f_\star]$,
\[
\big[\left(\SS_i, G_i\right), \f_i\big] \to \big[\left(S, G_\star\right), \f_\star \circ \wh \f_\star\big]=\big[\left(S, G_\star\right), \wh \f_\star\big] \quad \text{ in }\Tc(S).
\]
For any $\e>0$, and for all large enough $i$, we have~\eqref{Eq:dilatation} and
\[
\Ec_\Cc\big(\left[\left(\wh \f_\star\right)^\ast G_\star\right]\big) \le E_{\left(\wh \f_\star\right)^\ast G_\star}\left(h_i\right) \le (1+\e)E_{\f_i^\ast G_i}(h_i),
\]
and conclude that $\Ec_\Cc([(\wh \f_\star)^\ast G_\star]) \le R$ by letting $\e \to 0$, as required.
\end{proof}

\appendix

\section{The first and second variation formulas}

We show the first and second variation formulae for finite weighted graphs into Riemannian manifolds. This result is proved by Kotani and Sunada~\cite[Theorem~2.1]{KotaniSunadaStandard}; we give a proof for the sake of convenience.


Let $X=(V, E, m_E)$ be a finite weighted graph and $(M, G)$ a Riemannian manifold. A map $f: X\to M$ is called \emph{piecewise $C^k$} if $f_{e}:=f|_e: [0,1]\to M$ is $C^k$ for each edge $e\in E$. We assume $k\geq 2$ throughout this appendix. Consider a smooth family of piecewise $C^k$-maps $f_{s}: X \to (M, G)$ for $s\in (-\epsilon, \epsilon)$ with $f:=f_{0}$ for $\e>0$. Let $f_e(s, t):=f_{s,\,e} (t)$ and $T_e(s, t):=(\p/\p t)f_e(s, t)$ for each edge $e\in E$. We denote the variation vector fields by $(\p/\p s)f_e(s, t)=V_e(s, t)$. We denote the Levi--Civita connection and its curvature tensor of $G$ by $\nabla$ and $R$, respectively.


\begin{lemm}\label{Lem:first and second}
For any $s \in (-\e, \e)$, we have the first variation formula:
\begin{multline}\label{Eq:firstVF}
\frac{d}{d s} E_G(f_{s})\\
=-2\sum_{e\,\in\,E}m_E(e)G(V_e, T_e)(s,0)-\sum_{e\,\in\,E}m_E(e)\int_0^1 G\left(V_e, \nabla_{T_e}T_e\right)(s,t) dt,
\end{multline}
and the second variation formula:
\begin{multline}\label{Eq:secondVF}
\frac{d^2}{d s^2} E_G(f_{s})=-2\sum_{e\,\in\,E}m_E(e)G\left(\nabla_{V_e} V_e, T_e\right)(s,0)\\
+\sum_{e\,\in\,E}m_E(e) \int_0^1 \big\{-G\left(\nabla_{V_e}V_e, \nabla_{T_e}T_e\right)+\left\|\nabla_{T_e}V_e\right\|_G^2 -G\left(R\left(V_e, T_e\right)T_e, V_e\right) \big\}(s, t)dt.
\end{multline}
\end{lemm}

\begin{proof}
First, we see that
\begin{equation}\label{f1}
\begin{aligned}
\frac{d}{d s} E_G(f_{s}) &=\frac{1}{2}\sum_{e\,\in\,E}m_E(e)\int_0^1 \frac{\p}{\p s}G\left(T_e, T_e\right)\,dt =\sum_{e\,\in\,E}m_E(e)\int_0^1 G\left(\nabla_{V_e}T_e, T_e\right)\, dtr\\
&=\sum_{e\,\in\,E}m_E(e)\int_0^1G\left(\nabla_{T_e}V_e, T_e\right)\, dt \qquad  \left(\text{use }[V_e, T_e]=0\right)\\
&=\sum_{e\,\in\,E}m_E(e)\int_0^1\frac{\p}{\p t}G(V_e, T_e)- G\left(V_e, \nabla_{T_e}T_e\right)\, dt\\
&=\sum_{e\,\in\,E}m_E(e)\big[G\left(V_e, T_e\right)\big]_{t=0}^{t=1}-\sum_{e\,\in\,E}m_E(e)\int_0^1 G\left(V_e, \nabla_{T_e}T_e\right)\, dt.
\end{aligned}
\end{equation}
Since
\begin{equation}\label{Eq:ap1}
V_e(s, 1)=V_{\wbar e}(s, 0) \quad \text{and} \quad T_e(s,1)=- T_{\wbar e}(s,0)
\end{equation}
and $m_E(e)=m_E(\overline{e})$, we have that
\begin{align*}
\sum_{e\,\in\,E}m_E(e)G\left(V_e, T_e\right)(s,1) &=\sum_{e\,\in\,E}m_E(\overline{e})G\left( V_{\overline{e}}, -T_{\wbar e}\right)(s,0)\\
&=-\sum_{e\,\in\,E}m_E(e)G\left( V_e, T_e\right)(s,0),
\end{align*}
and thus,
\begin{equation}\label{f2}
\sum_{e\,\in\,E}m_E(e)\big[G\left(V_e, T_e\right)\big]_{t=0}^{t=1}=-2\sum_{e\,\in\,E}m_E(e)G\left(V_e, T_e\right)(s,0).
\end{equation}
Combining~\eqref{f2} with~\eqref{f1}, we obtain the first variation formula~\eqref{Eq:firstVF}.


Next, we have that
\begin{multline}\label{f3}
\frac{d^2}{d s^2} E_G(f_{s})\\
\begin{aligned}
&=\frac{1}{2}\sum_{e\,\in\,E}m_E(e)\int_0^1 \frac{\p^2}{\p s^2}G\left(T_e, T_e\right)\, dt\\
&=\sum_{e\,\in\,E}m_E(e)\int_0^1 G\left(\nabla_{V_e}\nabla_{V_e}T_e, T_e\right)+\left\|\nabla_{V_e}T_e\right\|^2\,dt \\
&=\sum_{e\,\in\,E}m_E(e)\int_0^1 G\left(\nabla_{V_e}\nabla_{T_e}V_e, T_e\right)+\left\|\nabla_{T_e}V_e\right\|^2\,dt \quad \left(\text{use}\; [V_e,T_e]=0\right)\\
&=\sum_{e\,\in\,E}m_E(e)\int_0^1 G\left(\nabla_{T_e}\nabla_{V_e}V_e, T_e\right)+G\left(R\left(V_e, T_e\right)V_e, T_e\right)+\left\|\nabla_{T_e}V_e\right\|^2\,dt.
\end{aligned}
\end{multline}
The first term becomes
\begin{multline}\label{f4}
\sum_{e\,\in\,E}m_E(e)\int_0^1 G\left(\nabla_{T_e}\nabla_{V_e}V_e, T_e\right)\, dt\\
\begin{aligned}
&=\sum_{e\,\in\,E}m_E(e)\int_0^1 \frac{\p}{\p t}G\left(\nabla_{V_e}V_e, T_e\right)-G\left(\nabla_{V_e}V_e, \nabla_{T_e}T_e\right)\, dt \\
&=-2\sum_{e\,\in\,E}m_E(e)G\left(\nabla_{V_e}V_e, T_e\right)(s,0)-\sum_{e\,\in\,E}m_E(e)\int_0^1G\left(\nabla_{V_e}V_e, \nabla_{T_e}T_e\right)\, dt,
\end{aligned}
\end{multline}
where we used~\eqref{f2} and $\nabla_{V_e}V_e(s,1)=\nabla_{V_{\wbar e}}V_{\wbar e}(s,0)$. Substituting~\eqref{f4} to~\eqref{f3} and using the fact $G(R(V_e, T_e)V_e, T_e)=-G(R(V_e, T_e)T_e, V_e)$, we obtain the second variation formula~\eqref{Eq:secondVF}.
\end{proof}


By the first variation formula~\eqref{Eq:firstVF}, a map $f$ is a critical point of the energy functional $E_G$ if and only if
\begin{equation}\label{harmmap}
\nabla_{T_e}T_e=0\quad \text{and}\quad \sum_{e\,\in\,E_x} m_E(e)T_e(x)=0
\end{equation}
for any edge $e\in E$ and $x\in V$. We call the map satisfying~\eqref{harmmap} a harmonic map. Moreover, according to the second variation formula~\eqref{Eq:secondVF}, we see the Hessian of $E_G$ for a harmonic map $h: X\to (M,G)$ is given by
\[
{\Hess}_{E_G}(V,W)_h=\sum_{e\,\in\,E} m_E(e)\int_0^1 \big\{G\left(\nabla_{T_e}V_e, \nabla_{T_e}W_e\right) -G\left(R\left(V_e, T_e\right)T_e, W_e\right)\big\}\,dt,
\]
for any two variation vector fields $V, W$ of $h$. From the expression, we see that ${\Hess}_{E_G}$ is non-negative definite if $(M,G)$ has non-positive sectional curvature. If furthermore, $(M,G)$ is a compact Riemannian manifold of negative sectional curvature and the image of the harmonic map $h: X\to M$ is not homotopic to a point nor a closed circle, then ${\Hess}_{E_G}$ is non-degenerate at $h$ (see the proof of Theorem~\ref{Thm:KSharmonic}$\MK$\eqref{theo2.1.3}).



\section{Smooth dependence of harmonic maps}\label{AppB}

Let $(M,G)$ be a compact smooth Riemannian manifold. For any integer $k\geq 1$, we denote the set of all $C^k$-Riemannian metrics on $M$ by $\Met^k(M)$. Note that the space $\Met^k(M)$ is an open subset of the Banach space consisting of $C^k$-symmetric $(0,2)$-tensors on $M$.

\begin{prop}\label{smooth dependence}
Fix any integer $k \ge 1$. Let $X$ be a finite weighted graph, $(M, G)$ be a closed Riemannian manifold of nonpositive sectional curvature and $h: X\to M$ be a harmonic map. If the Hessian $\Hess_{E_G}$ of the energy functional $E_G$ is non-degenerate at $h$, then there exists an open neighborhood $\Uc$ of $G$ in $\Met^{k+1}(G)$ and $C^k$-map $\wt h:\Uc \to C^{k+2}(X, M)$ such that $\wt h(G'):X \to M$ is a harmonic map for each $G' \in \Uc$ and $\wt h(G)=h$.
\end{prop}

There is a corresponding result where the domain and the target are Riemannian manifolds by Eells--Lemaire~\cite[Theorem~3.1]{EL} and Koiso~\cite[Theorem~4.7]{Koiso}. We shall prove by a simple application of the implicit function theorem between Banach spaces as in~\cite{EL}. Since we have not found the argument adapted to our setting, we describe the setup in a self-contained manner in the following.



For any piecewise $C^k$-map $f:X \to M$, let us denote the vector space consisting of $C^k$-vector fields along $f$ by $C^k(X, f^{-1}TM)$. We understand that every $v$ in $C^k(X, f^{-1}TM)$ is in $C^k$ on each interior $(0, 1)$ of edges and continuous on $X$. We define a norm on $C^k(X, f^{-1}TM)$ by
\[
\|v\|_k:=\sum_{0\,\leq\,i\,\leq\,k}\underset{e\,\in\,E_0}\sup\;\sup_{t\,\in\,[0,\, 1]}\left\|\nabla^i_{\p_t f_e} v_e\right\|_{G},
\]
where $E_0$ is the set of unoriented edges, for each edge $e\in E_0$, $v_e$ is the restriction of $v$ and $\nabla^i_{\p_t f_e}$ denotes the $i^{\rm th}$ covariant derivative along $f_e$ relative to the metric $G$. Then, $(C^k(X, f^{-1}TM), \|\cdot\|_k)$ becomes a Banach space.


Let $C^k(X, M):=\{f:X\to M: f\ \text{ is a piecewise $C^k$-map}\}$. Then, $C^k(X, M)$ is a Banach manifold modeled on the Banach space $(C^k(X, f^{-1}TM), \|\cdot \|_k)$. Namely, for any $f\in C^k(X, M)$ and an open neighborhood $B\subset C^k(X, f^{-1}TM)$ of $0$, we define the \emph{exponential map} ${\exp}_f:B\to C^k(X,M)$ by $ (\exp_fv)(x):=\exp_{f(x)}v_{x}, $ where $\exp_{f(x)}: T_{f(x)}M\to M$ denotes the exponential map of the Riemannian manifold $(M,G)$. Then, for an open neighborhood $B$, the map $\exp_f^{-1}: \exp_f(B)\to B$ gives a local chart around the point $f$.



We consider the subset of $C^k(X, M)$ consisting of piecewise $C^k$-maps satisfying the balanced condition:
\[
C^k_{bal}(X,M):=\left\{f\in C^k(X,M): \sum_{e\,\in\,E_x}m_E(e)\p_t f_e (x)=0 \quad \text{for any }x\in V\right\},
\]
where $\p_t f_e=(d/dt)f_e$. Note that the set $C^k_{bal}(X,M)$ is not empty since there always exists a harmonic map, which satisfies the balanced condition. For any $f \in C^k_{bal}(X,M)$, let us define the Banach space
\[
\Vc_{f,\,G}^k:=\!\left\{v\in C^k\left(X, f^{-1}TM\right)\,:\, \nabla^G_{v_x}\cdot \left(\sum_{e\,\in\,E_{x}}m_E(e) \p_t{f}_e(x)\right)=0 \quad \text{for any }\;x\in V\right\},
\]
as a closed subspace of $C^k(X, f^{-1}TM)$. Note that the condition in the definition $\Vc_{f,\,G}^k$ is equivalent to
\[
\sum_{e\,\in\,E_{x}}m_E(e) \nabla^G_{\p_t{f}_e}v_e(x)=0 \quad \text{for any }x\in V
\]
since $[v_e, \p_t{f}_e]=0$ for any $e \in E$.


\begin{lemm}
For any integer $k \ge 1$, $C^k_{bal}(X,M)$ is a submanifold of $C^k(X,M)$ and for any $C^k$-metric $G$ on $M$ and any $f \in C^k_{bal}(X,M)$, the tangent space of $C^k_{bal}(X,M)$ at $f$ is isomorphic to the Banach space $\Vc_{f,\,G}^k$.
\end{lemm}


\begin{proof}
Fix an arbitrary $f\in C^k_{bal}(X, M)$. Since $\Vc_{f,\,G}^k$ is a closed subspace of $C^k(X, f^{-1}TM)$, it becomes a Banach space with the induced norm $\|\cdot\|_k$. For a sufficiently small $v\in C^k(X, f^{-1}TM)$ with respect to $\|\cdot\|_k$, we set $f^v:=\exp_f\, v\in C^k(X, M)$. First, we show $f^v\in C^k_{bal}(X, M)$ if $v\in \Vc_{f,\,G}^k$. We take a geodesic normal coordinate $(U,\{x^1,\,\ldots,\,x^m\})$ of $(M,G)$ around a point $f(x)$ for $x\in V$, and define
\[
\exp: U\times \R^m\to M\text{ by } \exp\left(p, v^1\,\ldots,\,v^m\right)=\exp_{p}\left(\sum_{i=1}^mv^i\frac{\p}{\p x^i}\middle|_{p}\right).
\]
By definition of $\exp$ and $f^v$, we have
\[
\frac{d}{d t} {f}^v_e(x)=\frac{d}{dt} \exp\left(f_e(t), v_e(t)\right)\Big|_{t=0}=\left(d\exp\right)_{\left(f(x), v_x\right)}\left(\p_t f_e(x), \nabla^G_{\p_t{f}_e}v_e(x)\right).
\]
for each $e\in E_x$, and thus we obtain
\begin{equation}\label{dexp}
\sum_{e\,\in\,E_x}m_E(e)\frac{d}{d t} f^v_e(x)=\left(d{\exp}\right)_{\left(f(x), \,v_x\right)}\left(0, \sum_{e\,\in\,E_x}m_E(e)\nabla^G_{\p_t{f}_e}v_e(x)\right)
\end{equation}
since $f\in C^k_{bal}(X, M)$. Therefore, if $v\in \Vc_{f,\,G}^k$, then $f^v\in C^k_{bal}(X, M)$ as required.

Thus, we obtain a well-defined map
\[
\exp_f: \Uc\cap \Vc_{f,\,G}^k\to \exp_f(\Uc)\cap C^k_{bal}(X,M)
\]
for an open neighborhood $\Uc$ in $C^k(X, f^{-1}TM)$ around $0$, and this map is injective if $\Uc$ is small enough since it is a restriction of the exponential map $\exp_f$ to $\Vc_{f,\,G}^k$. We shall show the map is surjective, i.e., for any $f^v\in \exp_f(\Uc)\cap C^k_{bal}(X, M)$, the corresponding vector field $v:=\exp^{-1}_f(f^v)\in C^k(X, f^{-1}TM)$ is actually contained in $\Vc_{f, G}^k$. By~\eqref{dexp}, it is sufficient to show that $(d\exp)_{(x,\,v_x)}(0,w)=0$ implies $w=0$. Indeed, we have
\[
(d\exp)_{\left(f(x),\,v_x\right)}(0,w)=d\left(\exp_{f(x)}\right)_{v_x}(w),
\]
and hence, if $\Uc$ is small enough so that for any $v \in \Uc$ satisfies that $\|v\|_k<\epsilon$ for sufficiently small $\epsilon>0$ (for instance, we may take $\epsilon$ as ${\inj}(M)>0$), then $d({\exp}_x)_{v_x}$ is injective, and $w=0$; as required. This implies the lemma.
\end{proof}


We define the map
\begin{gather*}
\t: C_{bal}^{k+2}(X, M) \times \Met^{k+1}(M) \to \prod_{e\,\in\,E_0}\left(C^k\left([0, 1], f_e^{-1}TM\right)\right),
\\
\t(f, G):=\left(\nabla_{\p_t f_e}^G \p_t f_e\right)_{e\,\in\,E_0},
\end{gather*}
where $E_0$ is the set of unoriented edges of $X$; for each $e \in E_0$, we identify $e$ with $[0, 1]$ and denote by $C^k([0, 1], f_e^{-1}TM)$ the space of $C^k$-vector fields along $f_e:[0, 1] \to M$.

Let us fix a piecewise smooth harmonic map $h: X\to (M,G)$. Taking a small enough open neighborhood $U_{bal}\subset \Vc_{h,\,G}^{k+2}$ of $0$, and we identify ${\exp}_h(U_{bal})\subset C^{k+2}_{bal}\linebreak(X, M)$ and $U_{bal}$. For any $f \in U_{bal}$, we identify
\[
\prod_{e\,\in\,E_0}\left(C^k\left([0, 1], f_e^{-1}TM\right)\right) \quad \text{and} \quad \prod_{e\,\in\, E_0}\left(C^k\left([0, 1], h_e^{-1}TM\right)\right)
\]
via parallel transport relative to $\nabla^G$ along the curve $\gamma_f(t,x):=\exp_{h(x)}(tv(x))$ where $f=\exp_h\, v\in U_{bal}$.



We take any smooth curve $s \mapsto (f_s, G_s)$ in $U_{bal}\times \Met^{k+2}(M)$ through $(h, G)$ at $s=0$ for $s \in (-\e, \e)$ and for $\e>0$. Let us write $\t(f_s, G_s)=(\t_{s,\,e})_{e\,\in\,E_0}$ for $s \in (-\e, \e)$. Then, by the equations~\eqref{c4} and~\eqref{c41} given in the proof of Lemma~\ref{le3} (Note that the equations~\eqref{c4} and~\eqref{c41} hold for any Riemannian manifold as a target manifold), we see on each $e \in E_0$,
\begin{equation}\label{derivtau}
\frac{d\tau_{s,\,e}}{ds}\left|_{s=0}=\nabla_{T_e}\nabla_{T_e}V_e+R\left(V_e, T_e\right)T_e+\frac{d}{ds}\right|_{s=0} \nabla^{G_s}_{T_e}T_e,
\end{equation}
where $T_e:=\p_t h_e$ and $V_e:=(\partial_s f_{s, \,e})|_{s=0}$. Note that the final term in the right hand side depends only on $G_s$ and the initial condition $h$.

\begin{proof}[Proof of Proposition~\ref{smooth dependence}]
For a harmonic map $h:X \to (M, G)$, we take an open neighborhood $U_{bal}$ of $h$ in $\Vc_{h,\,G}^{k+2}$ and an open neighborhood $\Uc\subset \Met^{k+1}(M)$ of $G$. We regard $\tau$ as a $C^k$-map between Banach spaces:
\[
\tau: U_{bal}\times \Uc\to \prod_{e\,\in\,E_0}\left(C^k\left([0, 1], h_e^{-1}TM\right)\right).
\]
We shall show that the derivative of $\t$ in the $U_{bal}$-component at $o=(h, G)$
\[
d\tau|_{U_{bal},\,o}: \Vc_{h,\,G}^{k+2}\to \prod_{e\,\in\,E_0}\left(C^k\left([0, 1], h_e^{-1}TM\right)\right)
\]
is an isomorphism.

By~\eqref{derivtau}, we have that $d\tau_e|_{U_{bal},\,o}(V)=\nabla_{T_e}\nabla_{T_e}V_e+R(V_e, T_e)T_e$ on each $e \in E_0$ for any $V\in \Vc_{h,\, G}^{k+2}$ since the last term in~\eqref{derivtau} is independent of the deformation $f_s$ of $h$ relative to $V$. Thus, $d\tau|_{U_{bal},\,o}$ is a bounded linear map. First we show that $d\tau|_{U_{bal},\,o}$ is injective. Indeed, if $d\tau|_{U_{bal},\,o}(V)=0$, then taking the deformation $f_s$ relative to $V$ in $C^{k+2}_{bal}(X, M)$ for any $s\in (-\e, \e)$ for $\e>0$, we see that
\begin{align*}
0&=\sum_{e\,\in\,E}m_E(e)\int_0^1 G\left(-\nabla_{T_e}\nabla_{T_e}V_e-R\left(V_e, T_e\right)T_e, V_e\right)\, dt\\
&\quad+\sum_{e\,\in\,E}m_E(e)G_0\left(\nabla_{T_e}V_e, V_e\right)\Big|_{t=0}^{t=1}\\
&={\Hess}_{E_G}(V,V)_h,
\end{align*}
where we used Lemma~\ref{le2} in the first equality which follows from the assumption $f_s\in C^{k+2}_{bal}(X,M)$. The assumption in the statement implies that ${\Hess}_{E_G}$ is non-degenerate at $h$, we obtain $V=0$, and thus $d\tau|_{U_{bal},\,o}$ is injective.

Next we show that $d\tau|_{U_{bal},\,o}$ is surjective. For any given $W=(W_e)_{e \in E_0}$ where $W_e \in C^k([0, 1], h_e^{-1}TM)$ for each $e \in E_0$, we consider the second order ordinary differential equations
\begin{equation}\label{dtauw}
\nabla_{T_e}\nabla_{T_e}V_e+R(V_e, T_e)T_e=W_e
\end{equation}
on each $e \in E_0$ with the condition that $V$ is in $\Vc_{h,\,G}^{k+2}$, i.e., $V$ solves~\eqref{dtauw} on each interior $(0, 1)$ of edges $e$, and is continuous on $X$ satisfying the balanced condition. Let us consider the Hilbert space $H^{1, \,2}_{h,\,G}$ as the completion of $\Vc_{h,\,G}^{k+2}$ endowed with the inner product
\[
\left\langle V^1, V^2\right\rangle_{H_{h,\,G}^{1,\,2}}:=\sum_{e \,\in\,E}m_E(e)\int_0^1\Big\{G\left(\nabla_{T_e}V_e^1, \nabla_{T_e}V_e^2\right)-G\left(R\left(V_e^1, T_e\right)T_e, V_e^2\right)\Big\}\,dt,
\]
where this indeed defines the positive definite inner product since $(M, G)$ has nonpositive sectional curvature and ${\Hess}_{E_G}$ is non-degenerate at $h$. Note that every $V$ in $H^{1,\,2}_{h,\,G}$ is continuous on $X$. For any $W \in \prod_{e\,\in\,E_0}(C^k([0, 1], h_e^{-1}TM))$, we define a linear functional $L_W$ on $H^{1,\,2}_{h,\, G}$ by

\[
L_W(\f):=\sum_{e\,\in\,E}m_E(e)\int_0^1 G\left(W_e, \f_e\right)\,dt, \quad \text{for }\f \in H^{1,\,2}_{h,\,G}.
\]
Then, $L_W$ is a bounded linear functional on $H^{1, 2}_{h,\,G}$; in fact, the natural embedding map $H^{1,\,2}_{h,\,G} \to C^0(X, h^{-1}TM)$ is compact. Therefore the Riesz representation theorem implies that there exists a unique $V$ in $H^{1,\,2}_{h,\,G}$ such that $\langle V, \f\rangle_{H_{h,\,G}^{1,\,2}}=-L_W(\f)$ for any $\f$ in $H^{1,\,2}_{h,\,G}$. This implies that $V$ solves~\eqref{dtauw} on each interior $(0, 1)$ of edges $e$ and $V_e$ is in $C^{k+2}$ on $(0, 1)$ by taking smooth functions $\f_e$ whose supports are included in $(0, 1)$ for each $e \in E_0$. Then, integration by parts gives
\begin{multline*}
\sum_{e\,\in\,E}m_E(e)\int_0^1\big\{G\left(-\nabla_{T_e}\nabla_{T_e}V_e-R\left(V_e, T_e\right)T_e, \f_e\right)\big\}\,dt -2\sum_{x\,\in\,V}\sum_{e\,\in\,E_x}G\left(\nabla_{T_e}V_e, \f_e\right)\\
=-\sum_{e\,\in\,E}m_E(e)\int_0^1 G\left(W_e, \f_e\right)\,dt,
\end{multline*}
and $\sum_{x\,\in\,V}\sum_{e\,\in\,E_x}G(\nabla_{T_e}V_e, \f_e)=0$ holds for any $\f$ in $H^{1,\,2}_{h,\,G}$. Therefore $V$ satisfies the balanced condition and is in $V_{h,\,G}^{k+2}$. Hence $d\tau|_{U_{bal},\,o}$ is surjective.



Now, the implicit function theorem for Banach spaces implies that there exist an open neighborhood $\Uc'\subset \Uc$ of $G$ and a unique $C^k$-map $\wt h: \Uc'\to U_{bal}$ satisfying that
\[
\tau\left(\wt h\left(G'\right),G'\right)=0
\]
on $\Uc'$. Then, $\wt h(G): X\to M$ is a harmonic map relative to the metric $G'$ for any $G' \in \Uc'$ and $\wt h(G)=h$, and we complete the proof.
\end{proof}

\begin{thebibliography}{MMMMMM}
\bibitem[CdV91a]{ColindeVerdiereComment} Colin de Verdière, Yves Comment rendre géodésique une triangulation d’une surface ?, Enseign. Math., Volume 37 (1991) no. 3-4, pp. 201-212 \href{https://ahl.centre-mersenne.org/item/AHL_2021__4__1767_0/#ColindeVerdiereComment}{Publisher reference}.
\bibitem[CdV91b]{ColindeVerdiere} Colin de Verdière, Yves Un principe variationnel pour les empilements de cercles, Invent. Math., Volume 104 (1991) no. 3, pp. 655-669 \href{https://ahl.centre-mersenne.org/item/AHL_2021__4__1767_0/#ColindeVerdiere}{Publisher reference}.
\bibitem[EL81]{EL} Eells, James Jun.; Lemaire, Luc Deformations of metrics and associated harmonic maps, Proc. Indian Acad. Sci., Math. Sci., Volume 90 (1981) no. 1, pp. 33-45 \href{https://ahl.centre-mersenne.org/item/AHL_2021__4__1767_0/#EL}{Publisher reference}.
\bibitem[Eps66]{Epstein-isotopies} Epstein, David B. A. Curves von 2-manifolds and isotopies, Acta Math., Volume 115 (1966), pp. 83-107 \href{https://ahl.centre-mersenne.org/item/AHL_2021__4__1767_0/#Epstein-isotopies}{Publisher reference}.
\bibitem[FM11]{FarbMargalit} Farb, Benson; Margalit, Dan A Primer on Mapping Class Groups, Princeton Mathematical Series, 49, Princeton University Press, 2011 \href{https://ahl.centre-mersenne.org/item/AHL_2021__4__1767_0/#FarbMargalit}{Publisher reference}.
\bibitem[FN03]{FenchelNielsenBook} Fenchel, Werner; Nielsen, Jakob Discontinuous Groups of Isometries in the Hyperbolic Plane, De Gruyter Studies in Mathematics, 29, Walter de Gruyter, 2003 \href{https://ahl.centre-mersenne.org/item/AHL_2021__4__1767_0/#FenchelNielsenBook}{Publisher reference}.
\bibitem[Hat02]{Hatcher} Hatcher, Allen Algebraic Topology, Cambridge University Press, 2002 \href{https://ahl.centre-mersenne.org/item/AHL_2021__4__1767_0/#Hatcher}{Publisher reference}.
\bibitem[Kee74]{Keen} Keen, Linda Collars on Riemann surfaces, Discontinuous Groups and Riemann Surfaces (Proc. Conf., Univ. Maryland, College Park, Md., 1973) (Annals of Mathematics Studies), Volume 79, Princeton University Press (1974), pp. 263-268 \href{https://ahl.centre-mersenne.org/item/AHL_2021__4__1767_0/#Keen}{Publisher reference}.
\bibitem[Ker83]{KerckhoffNielsen} Kerckhoff, Steven P. The Nielsen realization problem, Ann. Math., Volume 117 (1983) no. 2, pp. 235-265 \href{https://ahl.centre-mersenne.org/item/AHL_2021__4__1767_0/#KerckhoffNielsen}{Publisher reference}.
\bibitem[Koi79]{Koiso} Koiso, Nohirito Variation of harmonic mapping caused by a deformation of Riemannian metric, Hokkaido Math. J., Volume 8 (1979) no. 2, pp. 199-213 \href{https://ahl.centre-mersenne.org/item/AHL_2021__4__1767_0/#Koiso}{Publisher reference}.
\bibitem[KS01]{KotaniSunadaStandard} Kotani, Motoko; Sunada, Toshikazu Standard realizations of crystal lattices via harmonic maps, Trans. Am. Math. Soc., Volume 353 (2001) no. 1, pp. 1-20 \href{https://ahl.centre-mersenne.org/item/AHL_2021__4__1767_0/#KotaniSunadaStandard}{Publisher reference}.
\bibitem[KWZ18]{Kim_et_al} Kim, Inkang; Wan, Xueyuan; Zhang, Genkai Plurisubharmonicity and geodesic convexity of energy function on Teichmüller space (2018) (\url{https://arxiv.org/abs/1809.00255v1}) \href{https://ahl.centre-mersenne.org/item/AHL_2021__4__1767_0/#Kim_et_al}{Publisher reference}.
\bibitem[Ste05]{Stephenson} Stephenson, Kenneth Introduction to Circle Packing: The Theory of Discrete Analytic Functions, Cambridge University Press, 2005 \href{https://ahl.centre-mersenne.org/item/AHL_2021__4__1767_0/#Stephenson}{Publisher reference}.
\bibitem[Sti92]{Stillwell} Stillwell, John C. Geometry of Surfaces, Universitext, Springer, 1992 \href{https://ahl.centre-mersenne.org/item/AHL_2021__4__1767_0/#Stillwell}{Publisher reference}.
\bibitem[Sun13]{SunadaTopological} Sunada, Toshikazu Topological Crystallography. With a View Towards Discrete Geometric Analysis, Surveys and Tutorials in the Applied Mathematical Sciences, 6, Springer, 2013 \href{https://ahl.centre-mersenne.org/item/AHL_2021__4__1767_0/#SunadaTopological}{Publisher reference}.
\bibitem[Thu78]{Thurston} Thurston, William P. The Geometry and Topology of Three-Manifolds (1978) (\url{http://www.math.unl.edu/~mbrittenham2/classwk/990s08/public/thurston.notes.pdf/6a.pdf}) \href{https://ahl.centre-mersenne.org/item/AHL_2021__4__1767_0/#Thurston}{Publisher reference}.
\bibitem[Tro92]{TrombaBook} Tromba, Anthony J. Teichmüller Theory in Riemannian Geometry: based on lecture notes by Jochen Denzler, Lectures in Mathematics, ETH Zürich, Birkhäuser, 1992 \href{https://ahl.centre-mersenne.org/item/AHL_2021__4__1767_0/#TrombaBook}{Publisher reference}.
\bibitem[Wol87]{WolpertNielsen} Wolpert, Scott A. Geodesic length functions and the Nielsen problem, J. Differ. Geom., Volume 25 (1987) no. 2, pp. 275-296 \href{https://ahl.centre-mersenne.org/item/AHL_2021__4__1767_0/#WolpertNielsen}{Publisher reference}.
\bibitem[Wol12]{Wolf} Wolf, Michael The Weil–Petersson Hessian of length of Teichmüller space, J. Differ. Geom., Volume 91 (2012) no. 1, pp. 129-169 \href{https://ahl.centre-mersenne.org/item/AHL_2021__4__1767_0/#Wolf}{Publisher reference}.
\bibitem[Yam99]{Yamada} Yamada, Sumio Weil–Petersson convexity of the energy functional on classical and universal Teichmüller spaces, J. Differ. Geom., Volume 51 (1999) no. 1, pp. 35-96 \href{https://ahl.centre-mersenne.org/item/AHL_2021__4__1767_0/#Yamada}{Publisher reference}.
\bibitem[Yam14]{Yamada2} Yamada, Sumio Local and global aspects of Weil–Petersson geometry, IRMA Lectures in Mathematics and Theoretical Physics, 19, European Mathematical Society (2014), pp. 43-111 \href{https://ahl.centre-mersenne.org/item/AHL_2021__4__1767_0/#Yamada2}{Publisher reference}.
\bibitem[Yam17]{YamadaSugaku} Yamada, Sumio On the Weil–Petersson convex geometry of Teichmüller space, Sugaku Expo., Volume 30 (2017) no. 2, pp. 159-186 \href{https://ahl.centre-mersenne.org/item/AHL_2021__4__1767_0/#YamadaSugaku}{Publisher reference}.
\end{thebibliography}
\end{document}
